% 2.5 章 · 傅里叶分析与 FFT
% 这一章要把"用计算机做傅里叶分析"拆成四件互不相同的数学事实：
%   1) DFT 是一个具体的 N×N 矩阵：正反变换、Parseval、位移定理都是它的代数性质；
%   2) Cooley-Tukey 把 O(N^2) 变成 O(N log N)，但下标约定（位反序）必须说清楚；
%   3) 循环卷积定理是恒等式，"等于线性卷积"是附加条件（零填充到 N_a+M-1）；
%   4) 加窗、补零、分辨率、泄漏这四件事各自的代价都可以量化，而且互不替代。
% 所有数字都来自本仓库 tests/2.5/ 的运行结果；断言附表由 tools/build_book.py 生成，
% \input 进来，不手抄。

\numlabchapter{傅里叶分析与 FFT}{2.5}

\section{这一章要解决什么}

给定长度为 $N$ 的有限序列 $x_0,\dots,x_{N-1}$，离散傅里叶变换（DFT）把时间下标换成频率下标：
\[
X_k=\sum_{n=0}^{N-1}x_n\omega_N^{kn},\qquad \omega_N=e^{-2\pi i/N},\qquad k=0,\dots,N-1 .
\]
这一层要回答的不是"DFT 是什么"（那是定义），而是四个工程问题：

\begin{enumerate}
  \item \textbf{怎么在 $O(N\log N)$ 里算出它，而且知道自己没算错？} 自写的 radix-2
        Cooley--Tukey（递归版与迭代版）与 $O(N^2)$ 定义式必须在**各种长度**上对拍：
        本章实测在 $n\in\{1,2,3,4,5,7,8,11,13,16,17,31,32,96,127,128,200,1000\}$ 上
        最坏相对误差 $1.46\times10^{-13}$（$n=1000$）。只测 2 的幂会漏掉整条 Bluestein 分支。
  \item \textbf{循环卷积定理什么时候等于"真的卷积"？} 定理说
        $\mathrm{DFT}(a\circledast b)=\mathrm{DFT}(a)\cdot\mathrm{DFT}(b)$ 恒成立；
        但它等于线性卷积的充要条件是 $N\ge N_a+N_b-1$。本章实测 $(N,M)=(16,8)$ 时
        卷绕误差 $0.251$（相对 $42.1\%$）且**只出现在前 $M-1=7$ 项**上——正好是你最会去看的头部；
        补零到 $23$ 之后降到 $1.22\times10^{-16}$。
  \item \textbf{窗、补零、分辨率、泄漏这四件事各自换了什么？} 它们经常被混为一谈：
        补零加密栅格（实测 $f_k$ 间距 $0.003906\to0.000489$）却不改分辨率；
        加窗压低旁瓣（矩形 $-13.26$ dB $\to$ Blackman $-58.11$ dB）却加宽主瓣
        （$0.884\to1.641$ 格）并压低相干增益（$1.0000\to0.4184$）。
  \item \textbf{"测到的频率"是什么？} 采样率 $f_s$ 之下的频率是等价类：实测 $3$ Hz 与 $17$ Hz
        在 $f_s=20$ 下的样本**逐位相同**（$\le10^{-12}$），幅度都读成 $1.0$。
\end{enumerate}

记号：$u=2^{-53}\approx1.11\times10^{-16}$ 是双精度单位舍入；$N$ 是序列长度；
$T=N\Delta t$ 是观测时长；$\Delta f_{\text{grid}}=1/(N\Delta t)$ 是频率栅格；
$D_N(x)=\dfrac{\sin(\pi x)}{N\sin(\pi x/N)}$ 是归一化 Dirichlet 核。
本章的"实测"数字全部取自 \texttt{tests/2.5/} 的断言记录。

\section{DFT 的代数：正反变换、Parseval、位移}

\subsection{反变换与正交性}

\begin{lemma}[根的单位性]
\label{lem:root-unity}
设 $\omega_N=e^{-2\pi i/N}$。则对任意整数 $m$，
\[
\sum_{n=0}^{N-1}\omega_N^{mn}=
\begin{cases}N,& m\equiv0\pmod N,\\ 0,& \text{否则}.\end{cases}
\]
\end{lemma}

\begin{proof}
若 $m\equiv0\pmod N$ 则每项为 $1$，和为 $N$。否则记 $r=\omega_N^m\ne1$，
由 $r^N=\omega_N^{mN}=e^{-2\pi im}=1$ 得
$\sum_{n=0}^{N-1}r^n=\frac{r^N-1}{r-1}=0$。
\end{proof}

\begin{theorem}[反变换公式与正交性]
\label{thm:inverse}
定义 $W_{kn}=\omega_N^{kn}$（$k,n=0,\dots,N-1$）。则 $W^{-1}=\frac1N\overline W$，即
\[
x_n=\frac1N\sum_{k=0}^{N-1}X_k\omega_N^{-kn},
\]
并且 $\sum_{k=0}^{N-1}|X_k|^2=N\sum_{n=0}^{N-1}|x_n|^2$（Parseval 恒等式）。
\end{theorem}

\begin{proof}
把 $X_k=\sum_m x_m\omega_N^{km}$ 代入右端：
\[
\frac1N\sum_{k=0}^{N-1}X_k\omega_N^{-kn}
=\frac1N\sum_m x_m\sum_{k=0}^{N-1}\omega_N^{k(m-n)}
=\frac1N\sum_m x_m\cdot N[m\equiv n]=x_n,
\]
最后一步用引理~\ref{lem:root-unity}。Parseval：由 $W^{-1}=N^{-1}\overline W$ 知
$W/\sqrt N$ 是酉矩阵，故 $\|Wx\|^2=N\|x\|^2$，即 $\sum_k|X_k|^2=N\sum_n|x_n|^2$。
\end{proof}

本章把这条件为回归测试：定义式正反变换的往返最坏误差 $3.83\times10^{-16}$，
正交性残差 $1.19\times10^{-16}$。实际实现里 $W$ 的元素写成
$\exp(-2\pi i\,(kn\bmod N)/N)$——**先做整数取模再取指数**：$kn$ 在 $n=1000$ 时可达 $10^6$，
直接算相位会先丢精度，取模以后 dft\_matrix 与 $\exp(-2\pi ikn/N)$ 的最大偏差是
$9.55\times10^{-16}$（$n=200$）。

\subsection{单边幅度谱的标定}

\begin{theorem}[幅度标定]
\label{thm:amplitude}
设 $x_n=A\cos(2\pi k_0 n/N+\varphi)$，$0<k_0<N/2$ 为整数。则单边幅度谱
\[
A_k=\frac{2|X_k|}{N}\ (0<k<N/2),\qquad A_0=\frac{|X_0|}{N},\qquad
A_{N/2}=\frac{|X_{N/2}|}{N}
\]
满足 $A_{k_0}=A$。若 $x_n\equiv c$ 则 $A_0=c$；若 $x_n=A(-1)^n$（Nyquist）则
$A_{N/2}=A$。
\end{theorem}

\begin{proof}
把余弦写成两个指数：$x_n=\frac A2e^{i\varphi}\omega_N^{-k_0n}+\frac A2e^{-i\varphi}\omega_N^{k_0n}$。
由引理~\ref{lem:root-unity}，第一项只在 $k=k_0$ 处给出 $N$，第二项只在 $k=N-k_0$ 处给出 $N$
（因为 $\omega_N^{k_0n}$ 的"频率"是 $N-k_0$）。于是 $X_{k_0}=\frac A2Ne^{i\varphi}$、
$X_{N-k_0}=\frac A2Ne^{-i\varphi}$，故 $|X_{k_0}|=AN/2$，乘 $2/N$ 得 $A$。
常数序列：$X_0=Nc$，其余为零，于是 $|X_0|/N=c$（**若乘 2 会读到 $2c$**）。
Nyquist：$x_n=A(-1)^n=A\omega_N^{(N/2)n}$，故 $X_{N/2}=NA$、$X_0=0$，
于是 $|X_{N/2}|/N=A$（**若乘 2 会读到 $2A$**）。
\end{proof}

三条路都要单独测，因为它们的错误方式不同：漏掉整体的 $2$ 让所有幅度小一半，
而 DC/Nyquist 也乘 $2$ 只让两端大一倍（中频正弦的读数完全正常）。
本章实测：12 组中频正弦 $+$ 纯直流（读 $2.5$，错法读 $5.0$）$+$ Nyquist（读 $1.75$，错法读 $3.5$），
最大读数偏差 $1.33\times10^{-15}$。

\begin{corollary}[单边 Parseval 的权重]
\label{cor:parseval-one-sided}
对实序列 $x$，令 $M=\lfloor N/2\rfloor$。则
\[
\Delta t\sum_{n=0}^{N-1}x_n^2=\frac{1}{N\Delta t}\Bigl(|X_0|^2+2\!\!\sum_{k=1}^{M-1}\!\!|X_k|^2+
[N\text{ 为偶}]\cdot|X_{N/2}|^2\Bigr).
\]
\end{corollary}

\begin{proof}
$x$ 实 $\Rightarrow X_{N-k}=\overline{X_k}$，故 $|X_{N-k}|=|X_k|$，把
定理~\ref{thm:inverse} 的 $\sum_{k=0}^{N-1}$ 折成单边时内部各出现两次、
$k=0$ 与（偶数 $N$ 时）$k=N/2$ 各出现一次。频率步长是
$\Delta f=1/(N\Delta t)$（**由原始长度决定**）。
\end{proof}

这里有两个独立的坑，必须分开断言。其一：把 $\Delta f$ 写成 $1/((N/2+1)\Delta t)$，
能量被放大 $N/(N/2+1)$ 倍——本章实测 $n=101$ 时比值 $1.9804$、Parseval 的相对误差
$0.980$：这个数字"不像粗心"（不是 $1$ 也不是 $10$），最容易被评为"数值误差"。
其二：内部谱线乘 $2$、DC/Nyquist 不乘；偶数 $N$ 有 Nyquist 项而奇数 $N$ 没有，
所以 $N$ 的奇偶都要测（实测 $n=64,\dots,1000$ 上相对误差 $\le6.9\times10^{-15}$）。

\subsection{位移定理}

\begin{theorem}[位移只改相位]
\label{thm:shift}
设 $y_n=x_{(n-p)\bmod N}$。则 $Y_k=\omega_N^{kp}X_k$，从而 $|Y_k|=|X_k|$；
对实序列，相位差 $\angle Y_k-\angle X_k=-2\pi kp/N$（模 $2\pi$）。
\end{theorem}

\begin{proof}
$Y_k=\sum_n x_{(n-p)\bmod N}\omega_N^{kn}$，令 $m=(n-p)\bmod N$，则 $n\equiv m+p$，
$\omega_N^{kn}=\omega_N^{k(m+p)}=\omega_N^{kp}\omega_N^{km}$，求和得 $\omega_N^{kp}X_k$。
\end{proof}

这条定理说明\textbf{幅度谱丢掉了时间原点}：本章实测位移 $5$ 点后幅度谱最大变化
$4.25\times10^{-16}$（应为 $0$），而相位斜坡的误差是 $2.27\times10^{-15}$。
要用 FFT 估时延，必须用相位或互相关，不能只看幅度。

\section{Cooley--Tukey：$O(N\log N)$ 与它附带的下标约定}

\subsection{递归 radix-2}

\begin{theorem}[DIT radix-2]
\label{thm:dit}
设 $N$ 为偶数，$x$ 的长度为 $N$。记偶下标子序列 $E_j=x_{2j}$、奇下标子序列 $O_j=x_{2j+1}$
（$j=0,\dots,N/2-1$），它们的 $N/2$ 点 DFT 为 $E_k,O_k$。则
\[
X_k=E_k+\omega_N^kO_k,\qquad X_{k+N/2}=E_k-\omega_N^kO_k,\qquad k=0,\dots,N/2-1 .
\]
\end{theorem}

\begin{proof}
把求和按 $n$ 的奇偶拆开：
\[
X_k=\sum_{j=0}^{N/2-1}x_{2j}\omega_N^{2jk}+\sum_{j=0}^{N/2-1}x_{2j+1}\omega_N^{(2j+1)k}
=\sum_j E_j(\omega_N^2)^{jk}+\omega_N^k\sum_j O_j(\omega_N^2)^{jk}.
\]
注意 $\omega_N^2=e^{-4\pi i/N}=e^{-2\pi i/(N/2)}=\omega_{N/2}$，两个和分别是
$E_k$ 与 $O_k$（在 $N/2$ 点 DFT 的意义下）。再由
$\omega_N^{k+N/2}=\omega_N^k\omega_N^{N/2}=-\omega_N^k$ 得第二式。
\end{proof}

\begin{corollary}[运算量]
\label{cor:complexity}
设 $T(N)$ 为 radix-2 的复数乘法次数。$T(N)=2T(N/2)+N/2$，$T(1)=0$，故
$T(N)=\tfrac N2\log_2N$；蝶形总数（每个蝶形一次复乘两次复加）也是 $\tfrac N2\log_2N$。
\end{corollary}

\begin{proof}
对 $T(N)=\tfrac N2+2T(N/2)$ 反复展开 $\log_2N$ 次：
$T(N)=\tfrac N2(1+1+\dots+1)$（共 $\log_2N$ 项）$=\tfrac N2\log_2N$。
\end{proof}

本章把这个账目**画进帧里并被断言**：$n=8/32/128$ 的级数各为 $3/5/7=\log_2N$，
蝶形数 $12/80/448=\tfrac N2\log_2N$，每级残差 $\le10^{-10}$（详见下文的"残差在量什么"）。

\subsection{迭代版与 DIF：位反序从哪来}

\begin{theorem}[DIF 的输出是位反序的]
\label{thm:dif}
设 $N=2^r$。考虑"频率抽取"分解：把输入的前后半相加/相减，
\[
u_m=x_m+x_{m+N/2},\qquad v_m=(x_m-x_{m+N/2})\,\omega_N^m,\qquad m=0,\dots,N/2-1 ,
\]
再对 $u$ 与 $v$ **各自递归**（两个长度 $N/2$ 的子问题，按下标连续分块）。
递归到底后，输出数组的第 $j$ 项是 $X_{\mathrm{rev}_r(j)}$，其中 $\mathrm{rev}_r$ 是 $r$ 位
二进制反转。对 $r=3$：
\[
\mathrm{rev}_3(0..7)=[0,4,2,6,1,5,3,7].
\]
\end{theorem}

\begin{proof}
对 $k$ 分奇偶，$k=2q$ 与 $k=2q+1$：
\[
X_{2q}=\sum_{n=0}^{N/2-1}\bigl(x_n+x_{n+N/2}\bigr)\omega_{N/2}^{nq},\qquad
X_{2q+1}=\sum_{n=0}^{N/2-1}\bigl(x_n-x_{n+N/2}\bigr)\omega_N^n\,\omega_{N/2}^{nq},
\]
即"偶频率"只依赖 $u$、"奇频率"只依赖 $v$，两者都是 $N/2$ 点 DFT。把 $k$ 的二进制最低位
当作这一级被决定的位，递归一次就固定一位；$r$ 级以后输出的下标顺序等于把
$k$ 的 $r$ 位二进制**反转**（最低位先被决定，于是跑到最高位）。
对 $r=3$：$j=0..7$ 对应 $k=\mathrm{rev}_3(j)$，即 $[0,4,2,6,1,5,3,7]$。
\end{proof}

本章的指纹检查用的是**复指数输入**而不是冲激：冲激的谱是常数，对任何置换都不敏感。
实测 $n=8$ 的指纹逐位等于 $\mathrm{rev}_3$，不做重排时 DIF 末行与 DFT 的相对误差是
$1.071$——这个 $O(1)$ 的"误差"其实是下标约定不同，不是精度问题。

\subsection{"每级残差"到底在量什么}

蝶形图那一页要报每一级的中间结果。本章的做法是：用两套**互相独立**的代码算同一批中间值
（一套是向量化的 DIF，一套是逐蝶形的朴素 DIF），逐级比较。实测残差 $n=4,8$ 时为 $0$，
$n=1024$ 时为 $2.85\times10^{-14}$。必须把这个数的含义说清楚：

\begin{quote}
它量的是\textbf{两份实现是否一致}，与 FFT 是否正确\textbf{无关}——两份一模一样写错的实现
残差同样是 $0$。数学正确性由另一条独立断言保证：与 $O(n^2)$ 定义式
\texttt{dft\_direct} 对拍，$n=1000$ 时相对误差 $1.46\times10^{-13}$。
\end{quote}

这是本章最容易被自己骗到的地方之一：一个漂亮的、随 $N$ 缓慢增长的"残差曲线"
看起来像是在展示"FFT 的数值稳定性"，其实它只是在展示"我抄了两遍同一份代码"。

\subsection{非 2 的幂：Bluestein}

\begin{theorem}[Bluestein 的长度条件]
\label{thm:bluestein}
设 $n$ 任意，取 $m\ge 2n-1$ 且 $m$ 为 $2$ 的幂。利用
$kj=\tfrac{k^2+j^2-(k-j)^2}{2}$，令
$a_j=x_j\omega_{2n}^{j^2}$、$b_j=\omega_{2n}^{-j^2}$（下标模 $m$），则
\[
X_k=\omega_{2n}^{-k^2}\sum_{j=0}^{n-1}a_jb_{k-j}
=\omega_{2n}^{-k^2}\,\bigl(a\circledast b\bigr)_k ,
\]
其中 $\circledast$ 是长度 $m$ 的循环卷积；$m\ge2n-1$ 保证它等于线性卷积
（见推论~\ref{cor:linear-eq-circular}），$m$ 是 $2$ 的幂保证用 radix-2 算。
\end{theorem}

\begin{proof}
$kj=\frac{k^2+j^2-(k-j)^2}{2}$ 是恒等式，故
$\omega_n^{kj}=\omega_{2n}^{k^2}\omega_{2n}^{j^2}\omega_{2n}^{-(k-j)^2}$。
代入 DFT 定义并提出与 $j$ 无关的因子即得第一式。$a$ 的支撑是 $[0,n-1]$、
$b$ 的支撑是 $[-(n-1),n-1]$，线性卷积的支撑长度是 $n+(2n-1)-1=3n-2$；
循环长度 $m\ge2n-1$ 时所需的下标 $k-j\in[-(n-1),n-1]$ 不会绕回来
（关键是用到的只是 $b$ 的那一段）。
\end{proof}

实测长度记账：$\texttt{bluestein\_length}(1,2,3,4096)=(1,4,8,8192)$，
全部 $\ge2n-1$ 且为 $2$ 的幂；$n=4096$ 走 Bluestein 与走 radix-2 的相对差
$2.43\times10^{-12}$（内部长度 $8192$）。

\begin{corollary}[Goertzel：只算一根谱线]
\label{cor:goertzel}
设 $\omega=2\pi k/N$，递推 $s_m=x_m+2\cos(\omega)s_{m-1}-s_{m-2}$（$s_{-1}=s_{-2}=0$）。
则
\[
X_k=e^{i\omega}s_{N-1}-s_{N-2}.
\]
\end{corollary}

\begin{proof}
先求冲激响应：$g_0=1$、$g_1=2\cos\omega$、$g_t=2\cos(\omega)g_{t-1}-g_{t-2}$，
解得 $g_t=\dfrac{\sin(\omega(t+1))}{\sin\omega}$（代入验证 $g_0,g_1$ 即可，
二阶线性递推由初值唯一确定）。由线性性
\[
s_m=\sum_{j=0}^{m}x_jg_{m-j}=\frac1{\sin\omega}\sum_{j=0}^{m}x_j\sin\bigl(\omega(m-j+1)\bigr).
\]
于是
\[
e^{i\omega}s_{N-1}-s_{N-2}
=\frac1{\sin\omega}\sum_{j=0}^{N-1}x_j\Bigl[e^{i\omega}\sin\bigl(\omega a\bigr)-\sin\bigl(\omega(a-1)\bigr)\Bigr],
\quad a=N-j .
\]
用 $\sin(\omega a)=\frac{e^{i\omega a}-e^{-i\omega a}}{2i}$ 展开括号：
\[
e^{i\omega}\sin(\omega a)-\sin\bigl(\omega(a-1)\bigr)
=\frac1{2i}\Bigl[e^{i\omega(a+1)}-e^{-i\omega(a-1)}-e^{i\omega(a-1)}+e^{-i\omega(a-1)}\Bigr]
=\frac1{2i}\Bigl[e^{i\omega(a+1)}-e^{i\omega(a-1)}\Bigr]
=e^{i\omega a}\sin\omega .
\]
代回：$e^{i\omega}s_{N-1}-s_{N-2}=\sum_j x_j e^{i\omega(N-j)}
=e^{i\omega N}\sum_j x_j e^{-i\omega j}=e^{2\pi ik}X_k=X_k$。
\end{proof}

\emph{这里 $e^{i\omega N}=e^{2\pi ik}=1$ 是整个推导的关键}：$\omega$ 必须是 $N$ 次单位根，
否则收尾式前面的因子不是 $1$。也正因为共轭方向只在
$k=0$ 与 $k=N/2$ 上不敏感（那里 $\omega\in\{0,\pi\}$，$e^{\pm i\omega}$ 相同），
本章的测试必须包含**随机信号 $+$ 一般 $k$**：用冲激或只用 $k=0,N/2$ 时，
把 $e^{i\omega}$ 写成 $e^{-i\omega}$ 也能"全绿"（实测 $k=1$ 时相对差 $1.29$）。

\section{卷积定理与它的适用条件}

\begin{theorem}[循环卷积定理]
\label{thm:circ-conv}
设 $a,b$ 长度均为 $N$，$(a\circledast b)_n=\sum_{m=0}^{N-1}a_{(n-m)\bmod N}b_m$。
则 $\mathrm{DFT}(a\circledast b)_k=A_kB_k$。
\end{theorem}

\begin{proof}
\[
\mathrm{DFT}(a\circledast b)_k=\sum_n\sum_m a_{(n-m)\bmod N}b_m\omega_N^{kn}
=\sum_m b_m\omega_N^{km}\sum_n a_{(n-m)\bmod N}\omega_N^{k(n-m)};
\]
内层求和令 $l=(n-m)\bmod N$，因 $\omega_N^{kN}=1$ 可去掉取模，得 $A_k$。
\end{proof}

\begin{corollary}[什么时候等于线性卷积]
\label{cor:linear-eq-circular}
设 $a$ 长 $N_a$、$b$ 长 $N_b$，把两者都补零到长度 $N$。若 $N\ge N_a+N_b-1$，则
循环卷积与线性卷积逐点相等（差别为 $0$）；否则差别只出现在下标
$0\le n\le N_b-2$（$M=N_b$）上，且尾部卷回头部。
\end{corollary}

\begin{proof}
线性卷积的支撑是 $[0,N_a+N_b-2]$，长度 $N_a+N_b-1$。当 $N\ge N_a+N_b-1$ 时
它整体落在 $[0,N-1]$ 内，取模下标 $(n-m)\bmod N$ 与真实下标 $n-m$ 相同，
两个求和逐项一致。否则下标 $n-m<0$ 的项被加上 $N$ 折回，只在
$n<m\le M-1$ 时发生，即 $0\le n\le M-2$。
\end{proof}

实测（$M=8$，高斯信号 $\times$ 均值滤波）：$(N,M)=(16,8)$ 时最大偏差 $0.251$
（相对 $42.1\%$）、差异位置数恰好 $7=M-1$；$N=24$ 时 $6.19\times10^{-2}$、
$N=32$ 时 $2.07\times10^{-2}$、$N=64$ 时 $2.63\times10^{-3}$——**它在变小，但不会变成 $0$**。
补零到 $N_a+N_b-1=23$ 后偏差 $1.22\times10^{-16}$（舍入地板）。
这就是"$N$ 取大一点"与"$N\ge N_a+N_b-1$"的区别：前者是量变，后者才是质变。

\section{窗、泄漏与分辨率}

\subsection{泄漏的机制}

\begin{theorem}[截断的谱就是 Dirichlet 核]
\label{thm:leakage}
设 $f_0=\kappa/T$（$\kappa$ 不必是整数），$x_n=e^{2\pi i f_0n\Delta t}$，$T=N\Delta t$。
则
\[
X_k=e^{i\pi(\kappa-k)(N-1)/N}\,D_N(\kappa-k),\qquad D_N(x)=\frac{\sin(\pi x)}{N\sin(\pi x/N)} .
\]
特别地：$\kappa$ 为整数时（整周期采样）除 $k=\kappa$ 外 $X_k=0$；
$|D_N|$ 的第一旁瓣峰值给出该偏移下的泄漏电平。
\end{theorem}

\begin{proof}
$X_k=\sum_{n=0}^{N-1}e^{2\pi i(\kappa-k)n/N}$ 是等比级数，公比
$r=e^{2\pi i(\kappa-k)/N}$：
\[
X_k=\frac{r^N-1}{r-1}
=e^{i\pi(\kappa-k)(N-1)/N}\cdot\frac{\sin(\pi(\kappa-k))}{\sin(\pi(\kappa-k)/N)},
\]
提出 $N$ 即得 $D_N$。$\kappa-k\in\mathbb Z$ 时分子 $\sin(\pi\cdot\text{整数})=0$，
故除 $\kappa=k$ 外全为 $0$（$k=\kappa$ 处极限为 $1$）。
\end{proof}

实测：$N=64$、频率恰好落在栅格上时次峰/主峰 $=0.000$（机器零）；
偏 $0.25$ 格时 $0.3265$（$-9.72$ dB）；偏 $0.5$ 格时 $0.7001$（$-3.10$ dB）。
所以\textbf{泄漏的原因是"信号周期与观测时长不成整数比"}，而不是"矩形窗不好"。
处理它只有三条路：改 $T$ 或 $f_s$ 让采样整周期、加窗（付主瓣变宽的代价）、
或者接受它并把它报出来。

\subsection{窗的代价：主瓣宽度、旁瓣、滚降、增益}

\begin{lemma}[端点光滑度决定滚降幂次]
\label{lem:rolloff}
设 $w:\mathbb R\to\mathbb R$ 有紧支撑且 $w\in C^m$（$w,w',\dots,w^{(m)}$ 在支撑端点外为零）。
则 $|\hat w(\omega)|\le C|\omega|^{-(m+1)}$（$|\omega|\to\infty$）。
\end{lemma}

\begin{proof}
对 $\hat w(\omega)=\int w(t)e^{-i\omega t}dt$ 分部积分 $m+1$ 次：
每次把 $-i\omega$ 提到分母并把导数移到 $w$ 上；由于
$w^{(j)}$（$j\le m$）在端点为零，边界项全部消失，剩下
$\hat w(\omega)=(-i\omega)^{-(m+1)}\int w^{(m+1)}(t)e^{-i\omega t}dt$，
而 $w^{(m+1)}$ 可积，故 $|\hat w|\le\|\hat{w^{(m+1)}}\|_\infty|\omega|^{-(m+1)}$。
\end{proof}

据此读本章的窗（离散版本用求和分部积分，幂次相同）：

\begin{center}
\begin{tabular}{lcccc}
窗 & 端点值 $w(0)$ & 端点一阶导 & 滚降 $p$ & 峰值旁瓣（实测/ Harris）\\
\hline
rect & $1$ & 跳跃 & $1$ & $-13.261$ / $-13.26$ dB\\
hann & $0$ & $0$ & $3$ & $-31.467$ / $-31.47$ dB\\
hamming & $0.08$ & 跳跃 & $1$ & $-42.675$ / $-42.68$ dB\\
blackman & $0$ & $0$ & $3$ & $-58.109$ / $-58.11$ dB\\
\end{tabular}
\end{center}

汉明窗是最好的反例：它把第一旁瓣压到 $-42.7$ dB（比 Hann 好 $11$ dB），
但端点不为零，于是 $p=1$——**压得低与滚降快是两件事**。
实测 $10\to1000$ 的包络比：矩形 $1.00\times10^{-2}$、Hann $1.00\times10^{-6}$。

\begin{theorem}[加窗的幅度代价 = 相干增益]
\label{thm:coherent-gain}
设 $x_n=A\cos(2\pi k_0n/N)$，$w$ 为窗。则加窗后单边幅度谱在 $k_0$ 处的读数
$A\cdot\mathrm{CG}(w)$，其中 $\mathrm{CG}(w)=\frac1N\sum_n w_n$；
能量 $\sum_n w_n^2 x_n^2$ 的期望则正比于 $\frac1N\sum_nw_n^2$。
\end{theorem}

\begin{proof}
加窗后 $X_{k_0}=\frac A2\sum_n w_n+\text{(镜像项)}$，第一项是窗的直流分量；
乘 $2/N$ 得 $A\cdot\frac1N\sum_nw_n$。能量的说法同理由
$\mathbb E[x_n^2]=A^2/2$ 与 $\sum_n w_n^2x_n^2\approx\frac{A^2}2\sum_nw_n^2$ 得到。
\end{proof}

实测：rect $1.0000$、hann $0.4980$、hamming $0.5382$、blackman $0.4184$；
每个窗的谱峰除以自己的 $\mathrm{CG}$ 后都回到 $2.0$（$n=256$，$A=2$）。
能量：rect $=N$、hann $=0.3746N$（连续极限 $3/8$）、blackman $=0.3043N$。
忘记除 $\mathrm{CG}$ 的效果很隐蔽：它不改旁瓣、不改频率位置，只让你报出去的幅度
系统性偏小一半。

\subsection{分辨率：Rayleigh 判据与它的边界}

\begin{theorem}[等幅双音的 Rayleigh 判据]
\label{thm:rayleigh}
设两个等幅单频信号频率相差 $\Delta\kappa/T$（以栅格为单位，$\kappa_1,\kappa_2$ 整数）。
用矩形窗截断后，其连续谱是两个 Dirichlet 核之和
$D_N(x-\kappa_1)+D_N(x-\kappa_2)$。当 $\kappa_2-\kappa_1=1$ 时，
两音中点的谱值**恰好为 $0$**；当间隔 $<1$ 格时两核在中点不再反号抵消，
两峰会合并进同一个主瓣。
\end{theorem}

\begin{proof}
由定理~\ref{thm:leakage}，单音的连续谱是 $D_N(x-\kappa)$。取中点
$x=(\kappa_1+\kappa_2)/2$，令 $d=(\kappa_2-\kappa_1)/2$，则两核在该点的值是
$D_N(-d)$ 与 $D_N(d)$。$D_N$ 是奇函数（分子 $\sin$ 奇、分母 $\sin$ 奇），故两者互为相反数，
和为 $0$：中点是严格的谷。$d<1/2$（间隔小于一格）时两核在中点同号（都不为零），
和不为零，谷消失。
\end{proof}

实测：整数格扫描下四个窗的最小可分辨间隔都是 $1.000/T$（= Rayleigh 判据），
$\Delta f_{\min}/$（$-3$ dB 主瓣宽度）落在 $0.5\sim1.8$ 之间（rect $1.12$、hann $0.69$、
hamming $0.77$、blackman $0.60$）。但**这张表只在整数格位置上成立**：

\begin{quote}
把两音放到半格（$2.25/T$ 与 $3.75/T$，间隔 $1.5/T$）上，矩形窗仍判"分得开"
（谷/峰 $0.251$），Hann 直接并成一个包。更早的一次实现用二分法去找阈值，
得到 $1.73/T$ 这个"看起来很像真值"的数——它是**假穿越**：谓词对间隔并不单调
（半格间隔会把两个音并进同一个主瓣），二分因此收敛到一个不存在的零点。
\end{quote}

所以"分辨率阈值"必须连同"在什么位置上量的"一起报；扫描代替二分，
或者先证明谓词单调。这也是本章把"补零不等于提高分辨率"写成机器判据的原因：
$\Delta f_{\text{grid}}$ 会被补零加密（$0.003906\to0.000489$，$pad=1\to8$），
而主瓣半高全宽的读数只从 $1.5625$ Hz 收敛到 $0.9766$ Hz（量化误差在消失），
它趋的是同一条连续主瓣，不是 $0$。

\section{补零：插值而不是分辨率}

\begin{theorem}[补零 = 三角插值]
\label{thm:zeropad}
设 $x$ 长 $N$，$x^{Z}$ 是把 $x$ 补零到长 $N'=pN$（$p\in\mathbb N$）。则
\[
X^{Z}_{pk}=X_k\qquad(k=0,\dots,N-1),
\]
即补零后的谱在**原栅格上逐点不变**；在其余点上它是
$\sum_n x_n e^{-2\pi i fn}$ 这条连续曲线的采样。
\end{theorem}

\begin{proof}
$X^Z_{pk}=\sum_{n=0}^{N'-1}x^Z_n e^{-2\pi ipkn/N'}=\sum_{n=0}^{N-1}x_n e^{-2\pi ikn/N}=X_k$，
因为 $n\ge N$ 的项为零，且 $e^{-2\pi ipkn/(pN)}=e^{-2\pi ikn/N}$。
\end{proof}

实测：$n=16..100$、$p=2,4,8$ 上 $\mathrm{fft}(\text{补零})[::p]$ 与 $\mathrm{fft}(x)$ 的
最坏相对差 $3.32\times10^{-14}$（即逐位相同，等于定理的机器版本）。
这条定理也解释了补零为什么**不能**把两个近邻音分开：它只在已有采样点之间插值，
不产生新信息。

\emph{顺带一个归一化陷阱}：单边幅度谱的 $2/N$ 里的 $N$ 必须是**原始长度**。
补零后若用补零后的长度归一化，每个峰都掉 $p$ 倍（实测 $p=2$ 时峰 $0.5$、$p=4$ 时 $0.25$），
而\textbf{谱的形状完全正常}——这不是"某个峰不对"，是整体比例不对，用眼睛很难发现。
本章的做法是显式传 \texttt{ref\_length}$=N$，并把"峰值不随 $p$ 改变"写成断言。

\section{采样：混叠是等价类}

\begin{theorem}[混叠]
\label{thm:aliasing}
以采样间隔 $\Delta t=1/f_s$ 采样 $x(t)=e^{2\pi ift}$，得到
$x_n=e^{2\pi ifn/f_s}$。则 $f$ 与 $f+kf_s$（$k\in\mathbb Z$）给出**完全相同**的样本；
在 $[0,f_s)$ 内可以取 $f_{\text{alias}}=|f-kf_s|$ 折回。
\end{theorem}

\begin{proof}
$e^{2\pi i(f+kf_s)n/f_s}=e^{2\pi ifn/f_s}\cdot e^{2\pi ikn}=e^{2\pi ifn/f_s}$。
\end{proof}

实测：$f_s=20$、$n=80$（$T=4$ s，两频率都是整周期）下 $3$ Hz 与 $17$ Hz 的样本
最大差 $\le10^{-12}$（应为 $0$），谱峰都在 $3/T$、幅度都读成 $1.0$；
$\texttt{alias\_frequency}(17,20)=3$ 而 $\texttt{alias\_frequency}(12,20)=8$
（$12$ 折到 $8$，不是 $3$——"高频都折到低频"是错的）。
"测到 $3$ Hz"因此只能说明真频率属于 $\{3,17,23,37,\dots\}$；要区分只能靠先验、
抗混叠滤波，或者改 $f_s$ 再看一次。

\section{五个可视化在看什么}

\begin{enumerate}
  \item \texttt{signal\_spectrum}（时频双视图）：$N$ 与 \texttt{pad} 分开控制。
        它要让人看到 $\Delta f_{\text{grid}}$ 随 \texttt{pad} 加密而
        $\Delta f_{\text{true}}=1/T$ 不动（帧 meta 里两个字段并排报），
        以及 $N$ 变小时基频与二次谐波如何并进同一个主瓣。
  \item \texttt{leak\_window}（泄漏与窗）：$N=64$、$M=64$ 的 dB 谱 + 四种窗的解析包络。
        理论线的锚点在窗自己的**最高旁瓣**上（rect $-13.26$ dB @ $fN=1.430$、
        hann $-31.47$ @ $2.360$、hamming $-42.68$ @ $4.500$、blackman $-58.11$ @ $3.545$），
        锚点与解析指标一致到 $3.5\times10^{-6}$ dB。
        注意汉明的锚点在 $fN=4.5$，**不是第一个零点 $fN=2$**：
        "第一个旁瓣"与"最高旁瓣"不是同一个位置。
  \item \texttt{butterfly\_stages}（分解与蝶形）：DIF 的级结构，每级一行、
        每级一个残差（残差的含义见 3.3 节）。$N=32$ 时 $5=\log_2 32$ 级、$80$ 个蝶形。
  \item \texttt{conv\_theorem}（卷积定理）：时域循环卷积、频域相乘、线性卷积三条曲线 + 两条差曲线。
        $N=16$、$M=8$ 时差曲线最大 $0.250781$（与 meta 逐位一致），
        补零到 $23$ 后 $1.11\times10^{-16}$。
  \item \texttt{resolution}（分辨率与不确定性）：两音 dB 谱、"能否分辨 vs 观测点数"的阶梯曲线，
        以及 Rayleigh 判据参考线。补零只把栅格加密 $8$ 倍而阈值漂移 $0.000/T$。
\end{enumerate}

% 本文件由 tools/build_book.py 从 dist/reports/assertions.json 生成，请勿手改。
\begin{longtable}{@{}p{0.06\textwidth}p{0.32\textwidth}p{0.22\textwidth}p{0.16\textwidth}@{}}
\caption{章节 2.5 数值断言（生成时间 2026-10-18T16:26:32-0400）}\\
\toprule
类别 & 断言 & 理论值 & 实测值 \\
\midrule
\endfirsthead
\toprule
类别 & 断言 & 理论值 & 实测值 \\
\midrule
\endhead
\bottomrule
\endlastfoot

B & signal\_\allowbreak{}spectrum 有 requiredAux error & $\le$ 1 & 1 \checkmark \\
B & signal\_\allowbreak{}spectrum 数组 x 全有限 & $\le$ 0 & 0 \checkmark \\
B & signal\_\allowbreak{}spectrum 数组 t\_\allowbreak{}axis 全有限 & $\le$ 0 & 0 \checkmark \\
B & signal\_\allowbreak{}spectrum 数组 spectrum\_\allowbreak{}db 全有限 & $\le$ 0 & 0 \checkmark \\
B & signal\_\allowbreak{}spectrum 数组 freq 全有限 & $\le$ 0 & 0 \checkmark \\
B & signal\_\allowbreak{}spectrum 数组 error 全有限 & $\le$ 0 & 0 \checkmark \\
B & leak\_\allowbreak{}window 有 requiredAux error & $\le$ 1 & 1 \checkmark \\
B & leak\_\allowbreak{}window 数组 window\_\allowbreak{}shapes 全有限 & $\le$ 0 & 0 \checkmark \\
B & leak\_\allowbreak{}window 数组 window\_\allowbreak{}shape 全有限 & $\le$ 0 & 0 \checkmark \\
B & leak\_\allowbreak{}window 数组 x\_\allowbreak{}windowed 全有限 & $\le$ 0 & 0 \checkmark \\
B & leak\_\allowbreak{}window 数组 spectrum\_\allowbreak{}db 全有限 & $\le$ 0 & 0 \checkmark \\
B & leak\_\allowbreak{}window 数组 freq 全有限 & $\le$ 0 & 0 \checkmark \\
B & leak\_\allowbreak{}window 数组 envelope 全有限 & $\le$ 0 & 0 \checkmark \\
B & leak\_\allowbreak{}window 数组 error 全有限 & $\le$ 0 & 0 \checkmark \\
B & butterfly\_\allowbreak{}stages 有 requiredAux residual & $\le$ 1 & 1 \checkmark \\
B & butterfly\_\allowbreak{}stages 数组 stage\_\allowbreak{}values\_\allowbreak{}real 全有限 & $\le$ 0 & 0 \checkmark \\
B & butterfly\_\allowbreak{}stages 数组 stage\_\allowbreak{}values\_\allowbreak{}imag 全有限 & $\le$ 0 & 0 \checkmark \\
B & butterfly\_\allowbreak{}stages 数组 current\_\allowbreak{}real 全有限 & $\le$ 0 & 0 \checkmark \\
B & butterfly\_\allowbreak{}stages 数组 current\_\allowbreak{}imag 全有限 & $\le$ 0 & 0 \checkmark \\
B & butterfly\_\allowbreak{}stages 数组 residual 全有限 & $\le$ 0 & 0 \checkmark \\
B & butterfly\_\allowbreak{}stages 数组 spectrum\_\allowbreak{}mag 全有限 & $\le$ 0 & 0 \checkmark \\
B & butterfly\_\allowbreak{}stages 数组 error 全有限 & $\le$ 0 & 0 \checkmark \\
B & butterfly\_\allowbreak{}stages 有 requiredAux residual & $\le$ 1 & 1 \checkmark \\
B & butterfly\_\allowbreak{}stages 数组 stage\_\allowbreak{}values\_\allowbreak{}real 全有限 & $\le$ 0 & 0 \checkmark \\
B & butterfly\_\allowbreak{}stages 数组 stage\_\allowbreak{}values\_\allowbreak{}imag 全有限 & $\le$ 0 & 0 \checkmark \\
B & butterfly\_\allowbreak{}stages 数组 current\_\allowbreak{}real 全有限 & $\le$ 0 & 0 \checkmark \\
B & butterfly\_\allowbreak{}stages 数组 current\_\allowbreak{}imag 全有限 & $\le$ 0 & 0 \checkmark \\
B & butterfly\_\allowbreak{}stages 数组 residual 全有限 & $\le$ 0 & 0 \checkmark \\
B & butterfly\_\allowbreak{}stages 数组 spectrum\_\allowbreak{}mag 全有限 & $\le$ 0 & 0 \checkmark \\
B & butterfly\_\allowbreak{}stages 数组 error 全有限 & $\le$ 0 & 0 \checkmark \\
B & conv\_\allowbreak{}theorem 有 requiredAux error & $\le$ 1 & 1 \checkmark \\
B & conv\_\allowbreak{}theorem 数组 signal 全有限 & $\le$ 0 & 0 \checkmark \\
B & conv\_\allowbreak{}theorem 数组 kernel 全有限 & $\le$ 0 & 0 \checkmark \\
B & conv\_\allowbreak{}theorem 数组 circular 全有限 & $\le$ 0 & 0 \checkmark \\
B & conv\_\allowbreak{}theorem 数组 linear 全有限 & $\le$ 0 & 0 \checkmark \\
B & conv\_\allowbreak{}theorem 数组 difference 全有限 & $\le$ 0 & 0 \checkmark \\
B & conv\_\allowbreak{}theorem 数组 safe\_\allowbreak{}difference 全有限 & $\le$ 0 & 0 \checkmark \\
B & conv\_\allowbreak{}theorem 数组 error 全有限 & $\le$ 0 & 0 \checkmark \\
B & resolution 有 requiredAux error & $\le$ 1 & 1 \checkmark \\
B & resolution 数组 spectrum\_\allowbreak{}db 全有限 & $\le$ 0 & 0 \checkmark \\
B & resolution 数组 freq 全有限 & $\le$ 0 & 0 \checkmark \\
B & resolution 数组 span\_\allowbreak{}ratio 全有限 & $\le$ 0 & 0 \checkmark \\
B & resolution 数组 span\_\allowbreak{}n 全有限 & $\le$ 0 & 0 \checkmark \\
B & resolution 数组 error 全有限 & $\le$ 0 & 0 \checkmark \\
B & 五个帧的数组与 aux & 全部有限且 requiredAux 齐备 & 5 \checkmark \\
B & pad=1 时栅格与真实分辨率一致 & 0.00390625 & 0.00390625 \checkmark \\
B & df = 1/(N*dt) & 0.00390625 & 0.00390625 \checkmark \\
B & 基频落在栅格上 & 768 & 768 \checkmark \\
B & 自写 FFT 与参照实现的最大偏差 & $\le$ 1.00000e-12 & $4.463\times 10^{-14}$ \checkmark \\
B & 默认参数下没有混叠 & $\le$ 0 & 0 \checkmark \\
B & pad=4 时栅格加密 4 倍 & 0.000976562 & 0.000976562 \checkmark \\
B & pad 不改变真实分辨率 df\_\allowbreak{}true & 0.00390625 & 0.00390625 \checkmark \\
B & pad 不改变主瓣宽度 & 0.00390625 & 0.00390625 \checkmark \\
B & signal\_\allowbreak{}spectrum 的 df 与 df\_\allowbreak{}true & df =\allowbreak{} 1/\allowbreak{}(\allowbreak{}N*dt) 随 pad 加密，df\_true =\allowbreak{} 1/\allowbreak{}T 不变 & 0.25 \checkmark \\
B & rect 帧里的 Harris 表值一致 & $\le$ 0.005 & 0 \checkmark \\
B & rect 帧里实测旁瓣与表值一致 & $\le$ 0.6 & 0.00145913 \checkmark \\
B & rect 帧里的相干增益 = 内核值 & 1 & 1 \checkmark \\
B & hann 帧里的 Harris 表值一致 & $\le$ 0.005 & 0 \checkmark \\
B & hann 帧里实测旁瓣与表值一致 & $\le$ 0.6 & 0.00269145 \checkmark \\
B & hann 帧里的相干增益 = 内核值 & 0.492188 & 0.492188 \checkmark \\
B & hamming 帧里的 Harris 表值一致 & $\le$ 0.005 & 0 \checkmark \\
B & hamming 帧里实测旁瓣与表值一致 & $\le$ 0.6 & 0.00496879 \checkmark \\
B & hamming 帧里的相干增益 = 内核值 & 0.532813 & 0.532813 \checkmark \\
B & blackman 帧里的 Harris 表值一致 & $\le$ 0.005 & 0 \checkmark \\
B & blackman 帧里实测旁瓣与表值一致 & $\le$ 0.6 & 0.00115825 \checkmark \\
B & blackman 帧里的相干增益 = 内核值 & 0.413438 & 0.413438 \checkmark \\
B & offset=0.0 时解析窗谱与 |FFT| 的偏差 $\le$ 5\% & $\le$ 0.05 & 0.0132386 \checkmark \\
B & offset=0.25 时解析窗谱与 |FFT| 的偏差 $\le$ 5\% & $\le$ 0.05 & 0.0132386 \checkmark \\
B & offset=0.5 时解析窗谱与 |FFT| 的偏差 $\le$ 5\% & $\le$ 0.05 & 0.0132386 \checkmark \\
B & rect 理论线的锚点 = 最高旁瓣 & $\le$ 0.01 & $3.525\times 10^{-6}$ \checkmark \\
B & hann 理论线的锚点 = 最高旁瓣 & $\le$ 0.01 & 0.000181148 \checkmark \\
B & hamming 理论线的锚点 = 最高旁瓣 & $\le$ 0.01 & 0.000231284 \checkmark \\
B & blackman 理论线的锚点 = 最高旁瓣 & $\le$ 0.01 & 0.000166771 \checkmark \\
B & 汉明的最高旁瓣在 fN $\approx$ 4.5，不是第一个零点 fN = 2 & $\le$ 4.5 & 4 \checkmark \\
B & 四个窗的帧内旁瓣与表值的最大差 & 不超过 0.6 dB & 0.00496879 \checkmark \\
B & n=8 级数 = log2(N) & 3 & 3 \checkmark \\
B & n=8 蝶形数 = (N/2)*log2(N) & 12 & 12 \checkmark \\
B & n=8 每级残差贴地板 & $\le$ 1.00000e-10 & 0 \checkmark \\
B & n=8 末级结果与参照谱一致 & $\le$ 1.00000e-12 & $2.890\times 10^{-16}$ \checkmark \\
B & n=32 级数 = log2(N) & 5 & 5 \checkmark \\
B & n=32 蝶形数 = (N/2)*log2(N) & 80 & 80 \checkmark \\
B & n=32 每级残差贴地板 & $\le$ 1.00000e-10 & $2.064\times 10^{-16}$ \checkmark \\
B & n=32 末级结果与参照谱一致 & $\le$ 1.00000e-12 & $2.398\times 10^{-15}$ \checkmark \\
B & n=128 级数 = log2(N) & 7 & 7 \checkmark \\
B & n=128 蝶形数 = (N/2)*log2(N) & 448 & 448 \checkmark \\
B & n=128 每级残差贴地板 & $\le$ 1.00000e-10 & $2.879\times 10^{-15}$ \checkmark \\
B & n=128 末级结果与参照谱一致 & $\le$ 1.00000e-12 & $1.124\times 10^{-14}$ \checkmark \\
B & 余弦 k=3 的谱峰 = N/2 = 16 & 16 & 16 \checkmark \\
B & 蝶形动画的 stage 随时间推进 & t 从 0 到 tMax 时 stage 单调不减 & 5 \checkmark \\
B & need = N+M-1 = 23 & 23 & 23 \checkmark \\
B & 补零后误差贴舍入地板 & $\le$ 1.00000e-15 & $1.110\times 10^{-16}$ \checkmark \\
B & 定理本身的残差贴地板 & $\le$ 1.00000e-12 & $1.145\times 10^{-16}$ \checkmark \\
B & 频域相乘与直接卷积一致 & $\le$ 1.00000e-12 & $1.110\times 10^{-16}$ \checkmark \\
B & 卷绕误差是 O(1) 量级而不是小量 & $\le$ 1 & 0.250781 \checkmark \\
B & N 增大后卷绕误差单调不增 & $\le$ 0.250781 & 0.00263061 \checkmark \\
B & 卷绕误差随 N 的变化 & N=\allowbreak{}16 时 O(\allowbreak{}1)，N=\allowbreak{}64 时已降到 1e-3 以下 & 95.3319 \checkmark \\
B & Rayleigh 判据 = 1/T & 1 & 1 \checkmark \\
B & sep\_\allowbreak{}measured = 内核的最小可分辨间隔(rect) & 1 & 1 \checkmark \\
B & sep\_\allowbreak{}over\_\allowbreak{}mainlobe = 阈值 / 主瓣宽度 & 1.1236 & 1.1236 \checkmark \\
B & pad=1 时栅格 = 真实分辨率 & 0.00390625 & 0.00390625 \checkmark \\
B & Hann 的阈值不小于矩形窗 & $\le$ 1 & 1 \checkmark \\
B & Hann 的主瓣比矩形宽 & $\le$ 1.44 & 0.89 \checkmark \\
B & pad=8 时栅格加密 8 倍 & 0.000488281 & 0.000488281 \checkmark \\
B & 补零与换窗对阈值的影响 & pad 只动栅格；阈值的变化远小于换窗带来的变化 & 0 \checkmark \\
B & dB 谱的峰落在 meta 报的 f0 上 & 3 & 3 \checkmark \\
B & signal\_\allowbreak{}spectrum 的 error 曲线贴地板 & $\le$ 1.00000e-12 & $4.463\times 10^{-14}$ \checkmark \\
B & 残差数组的最大值 = meta 的 max\_\allowbreak{}level\_\allowbreak{}residual & $2.064\times 10^{-16}$ & $2.064\times 10^{-16}$ \checkmark \\
B & 每级残差都不超过 1e-10 & $\le$ 1.00000e-10 & $2.064\times 10^{-16}$ \checkmark \\
B & difference 数组的最大值 = meta 的 max\_\allowbreak{}wrap & 0.250781 & 0.250781 \checkmark \\
B & safe\_\allowbreak{}difference 贴舍入地板 & $\le$ 1.00000e-15 & $1.110\times 10^{-16}$ \checkmark \\
B & 数组与 meta 的互证 & 最大值类读数必须逐位一致 & 0.250781 \checkmark \\
B & n=1 自研 FFT vs 定义式 DFT & $\le$ 7.10543e-15 & 0 \checkmark \\
B & n=1 自研 FFT vs pocketfft & $\le$ 7.10543e-15 & 0 \checkmark \\
B & n=2 自研 FFT vs 定义式 DFT & $\le$ 1.42109e-14 & $4.771\times 10^{-17}$ \checkmark \\
B & n=2 自研 FFT vs pocketfft & $\le$ 1.42109e-14 & 0 \checkmark \\
B & n=3 自研 FFT vs 定义式 DFT & $\le$ 2.13163e-14 & $3.419\times 10^{-16}$ \checkmark \\
B & n=3 自研 FFT vs pocketfft & $\le$ 2.13163e-14 & $3.516\times 10^{-16}$ \checkmark \\
B & n=4 自研 FFT vs 定义式 DFT & $\le$ 2.84217e-14 & $5.249\times 10^{-17}$ \checkmark \\
B & n=4 自研 FFT vs pocketfft & $\le$ 2.84217e-14 & $8.398\times 10^{-18}$ \checkmark \\
B & n=5 自研 FFT vs 定义式 DFT & $\le$ 3.55271e-14 & $3.447\times 10^{-16}$ \checkmark \\
B & n=5 自研 FFT vs pocketfft & $\le$ 3.55271e-14 & $5.407\times 10^{-16}$ \checkmark \\
B & n=7 自研 FFT vs 定义式 DFT & $\le$ 4.97380e-14 & $6.793\times 10^{-16}$ \checkmark \\
B & n=7 自研 FFT vs pocketfft & $\le$ 4.97380e-14 & $8.932\times 10^{-16}$ \checkmark \\
B & n=8 自研 FFT vs 定义式 DFT & $\le$ 5.68434e-14 & $2.687\times 10^{-16}$ \checkmark \\
B & n=8 自研 FFT vs pocketfft & $\le$ 5.68434e-14 & $1.202\times 10^{-16}$ \checkmark \\
B & n=11 自研 FFT vs 定义式 DFT & $\le$ 7.81597e-14 & $1.284\times 10^{-15}$ \checkmark \\
B & n=11 自研 FFT vs pocketfft & $\le$ 7.81597e-14 & $1.164\times 10^{-15}$ \checkmark \\
B & n=13 自研 FFT vs 定义式 DFT & $\le$ 9.23706e-14 & $1.433\times 10^{-15}$ \checkmark \\
B & n=13 自研 FFT vs pocketfft & $\le$ 9.23706e-14 & $1.474\times 10^{-15}$ \checkmark \\
B & n=16 自研 FFT vs 定义式 DFT & $\le$ 1.13687e-13 & $1.631\times 10^{-16}$ \checkmark \\
B & n=16 自研 FFT vs pocketfft & $\le$ 1.13687e-13 & $1.079\times 10^{-16}$ \checkmark \\
B & n=17 自研 FFT vs 定义式 DFT & $\le$ 1.20792e-13 & $1.740\times 10^{-15}$ \checkmark \\
B & n=17 自研 FFT vs pocketfft & $\le$ 1.20792e-13 & $1.754\times 10^{-15}$ \checkmark \\
B & n=31 自研 FFT vs 定义式 DFT & $\le$ 2.20268e-13 & $1.329\times 10^{-14}$ \checkmark \\
B & n=31 自研 FFT vs pocketfft & $\le$ 2.20268e-13 & $1.364\times 10^{-14}$ \checkmark \\
B & n=32 自研 FFT vs 定义式 DFT & $\le$ 2.27374e-13 & $2.342\times 10^{-16}$ \checkmark \\
B & n=32 自研 FFT vs pocketfft & $\le$ 2.27374e-13 & $1.171\times 10^{-16}$ \checkmark \\
B & n=96 自研 FFT vs 定义式 DFT & $\le$ 6.82121e-13 & $3.083\times 10^{-14}$ \checkmark \\
B & n=96 自研 FFT vs pocketfft & $\le$ 6.82121e-13 & $3.075\times 10^{-14}$ \checkmark \\
B & n=127 自研 FFT vs 定义式 DFT & $\le$ 9.02389e-13 & $4.728\times 10^{-14}$ \checkmark \\
B & n=127 自研 FFT vs pocketfft & $\le$ 9.02389e-13 & $4.766\times 10^{-14}$ \checkmark \\
B & n=128 自研 FFT vs 定义式 DFT & $\le$ 9.09495e-13 & $3.150\times 10^{-16}$ \checkmark \\
B & n=128 自研 FFT vs pocketfft & $\le$ 9.09495e-13 & $3.247\times 10^{-16}$ \checkmark \\
B & n=200 自研 FFT vs 定义式 DFT & $\le$ 1.42109e-12 & $6.313\times 10^{-15}$ \checkmark \\
B & n=200 自研 FFT vs pocketfft & $\le$ 1.42109e-12 & $6.271\times 10^{-15}$ \checkmark \\
B & n=1000 自研 FFT vs 定义式 DFT & $\le$ 7.10543e-12 & $1.456\times 10^{-13}$ \checkmark \\
B & n=1000 自研 FFT vs pocketfft & $\le$ 7.10543e-12 & $1.458\times 10^{-13}$ \checkmark \\
A & 最坏相对误差 & c$\cdot$n$\cdot$u & $1.456\times 10^{-13}$ \checkmark \\
B & n=1 往返 & $\le$ 7.10543e-15 & 0 \checkmark \\
B & n=2 往返 & $\le$ 1.42109e-14 & $1.420\times 10^{-16}$ \checkmark \\
B & n=3 往返 & $\le$ 2.13163e-14 & $4.714\times 10^{-16}$ \checkmark \\
B & n=4 往返 & $\le$ 2.84217e-14 & $1.267\times 10^{-16}$ \checkmark \\
B & n=5 往返 & $\le$ 3.55271e-14 & $2.123\times 10^{-16}$ \checkmark \\
B & n=7 往返 & $\le$ 4.97380e-14 & $7.078\times 10^{-16}$ \checkmark \\
B & n=8 往返 & $\le$ 5.68434e-14 & $1.764\times 10^{-16}$ \checkmark \\
B & n=11 往返 & $\le$ 7.81597e-14 & $8.984\times 10^{-16}$ \checkmark \\
B & n=13 往返 & $\le$ 9.23706e-14 & $1.555\times 10^{-15}$ \checkmark \\
B & n=16 往返 & $\le$ 1.13687e-13 & $3.039\times 10^{-16}$ \checkmark \\
B & n=17 往返 & $\le$ 1.20792e-13 & $2.117\times 10^{-15}$ \checkmark \\
B & n=31 往返 & $\le$ 2.20268e-13 & $4.623\times 10^{-15}$ \checkmark \\
B & n=32 往返 & $\le$ 2.27374e-13 & $2.448\times 10^{-16}$ \checkmark \\
B & n=96 往返 & $\le$ 6.82121e-13 & $1.624\times 10^{-14}$ \checkmark \\
B & n=127 往返 & $\le$ 9.02389e-13 & $1.906\times 10^{-14}$ \checkmark \\
B & n=128 往返 & $\le$ 9.09495e-13 & $3.683\times 10^{-16}$ \checkmark \\
B & n=200 往返 & $\le$ 1.42109e-12 & $1.937\times 10^{-14}$ \checkmark \\
B & n=1000 往返 & $\le$ 7.10543e-12 & $1.619\times 10^{-13}$ \checkmark \\
A & 忘记 1/N 时的往返放大倍数 & 应当恰好是 n & 1000 \checkmark \\
B & n=7 dft\_\allowbreak{}matrix 相位精度 & $\le$ 6.21725e-15 & 0 \checkmark \\
B & n=13 dft\_\allowbreak{}matrix 相位精度 & $\le$ 1.15463e-14 & $8.951\times 10^{-16}$ \checkmark \\
B & n=17 dft\_\allowbreak{}matrix 相位精度 & $\le$ 1.50990e-14 & $4.003\times 10^{-16}$ \checkmark \\
B & n=31 dft\_\allowbreak{}matrix 相位精度 & $\le$ 2.75335e-14 & $9.421\times 10^{-16}$ \checkmark \\
B & n=96 dft\_\allowbreak{}matrix 相位精度 & $\le$ 8.52651e-14 & $9.155\times 10^{-16}$ \checkmark \\
B & n=127 dft\_\allowbreak{}matrix 相位精度 & $\le$ 1.12799e-13 & $8.951\times 10^{-16}$ \checkmark \\
B & n=200 dft\_\allowbreak{}matrix 相位精度 & $\le$ 1.77636e-13 & $9.550\times 10^{-16}$ \checkmark \\
B & n=1000 dft\_\allowbreak{}matrix 相位精度 & $\le$ 8.88178e-13 & $9.550\times 10^{-16}$ \checkmark \\
A & dft\_\allowbreak{}matrix 最坏逐点误差 & O(\allowbreak{}n$\cdot$u) & $9.550\times 10^{-16}$ \checkmark \\
B & n=2 递归 vs 迭代 & $\le$ 2.84217e-14 & 0 \checkmark \\
B & n=2 递归 vs Bluestein & $\le$ 2.84217e-14 & 0 \checkmark \\
B & n=4 递归 vs 迭代 & $\le$ 5.68434e-14 & 0 \checkmark \\
B & n=4 递归 vs Bluestein & $\le$ 5.68434e-14 & 0 \checkmark \\
B & n=8 递归 vs 迭代 & $\le$ 1.13687e-13 & 0 \checkmark \\
B & n=8 递归 vs Bluestein & $\le$ 1.13687e-13 & 0 \checkmark \\
B & n=16 递归 vs 迭代 & $\le$ 2.27374e-13 & 0 \checkmark \\
B & n=16 递归 vs Bluestein & $\le$ 2.27374e-13 & 0 \checkmark \\
B & n=32 递归 vs 迭代 & $\le$ 4.54747e-13 & 0 \checkmark \\
B & n=32 递归 vs Bluestein & $\le$ 4.54747e-13 & 0 \checkmark \\
B & n=64 递归 vs 迭代 & $\le$ 9.09495e-13 & 0 \checkmark \\
B & n=64 递归 vs Bluestein & $\le$ 9.09495e-13 & 0 \checkmark \\
B & n=128 递归 vs 迭代 & $\le$ 1.81899e-12 & 0 \checkmark \\
B & n=128 递归 vs Bluestein & $\le$ 1.81899e-12 & 0 \checkmark \\
B & n=256 递归 vs 迭代 & $\le$ 3.63798e-12 & 0 \checkmark \\
B & n=256 递归 vs Bluestein & $\le$ 3.63798e-12 & 0 \checkmark \\
B & n=3 Bluestein vs 定义式 & $\le$ 4.26326e-14 & $2.698\times 10^{-16}$ \checkmark \\
B & n=5 Bluestein vs 定义式 & $\le$ 7.10543e-14 & $3.578\times 10^{-16}$ \checkmark \\
B & n=7 Bluestein vs 定义式 & $\le$ 9.94760e-14 & $5.891\times 10^{-16}$ \checkmark \\
B & n=11 Bluestein vs 定义式 & $\le$ 1.56319e-13 & $1.237\times 10^{-15}$ \checkmark \\
B & n=13 Bluestein vs 定义式 & $\le$ 1.84741e-13 & $1.221\times 10^{-15}$ \checkmark \\
B & n=17 Bluestein vs 定义式 & $\le$ 2.41585e-13 & $1.709\times 10^{-15}$ \checkmark \\
B & n=31 Bluestein vs 定义式 & $\le$ 4.40536e-13 & $1.356\times 10^{-14}$ \checkmark \\
B & n=100 Bluestein vs 定义式 & $\le$ 1.42109e-12 & $5.988\times 10^{-15}$ \checkmark \\
A & 三种实现的最大相对差 & 同一数学对象的不同算法 & $1.356\times 10^{-14}$ \checkmark \\
B & n=4 位反序重排后 vs DFT & $\le$ 1.13687e-13 & $5.763\times 10^{-17}$ \checkmark \\
B & n=8 位反序重排后 vs DFT & $\le$ 2.27374e-13 & $1.740\times 10^{-16}$ \checkmark \\
B & n=16 位反序重排后 vs DFT & $\le$ 4.54747e-13 & $1.294\times 10^{-16}$ \checkmark \\
B & n=32 位反序重排后 vs DFT & $\le$ 9.09495e-13 & $1.912\times 10^{-16}$ \checkmark \\
B & n=64 位反序重排后 vs DFT & $\le$ 1.81899e-12 & $1.969\times 10^{-16}$ \checkmark \\
B & n=128 位反序重排后 vs DFT & $\le$ 3.63798e-12 & $3.632\times 10^{-16}$ \checkmark \\
B & 重排后最坏相对误差 & $\le$ 3.63798e-12 & $3.632\times 10^{-16}$ \checkmark \\
A & 不重排时的相对误差 & 位反序排列与自然序的差别是 O(\allowbreak{}1) & 1.07082 \checkmark \\
B & n=4 指纹是置换 & $\le$ 0 & 0 \checkmark \\
B & n=4 指纹 = 位反序 & $\le$ 0 & 0 \checkmark \\
B & n=8 指纹是置换 & $\le$ 0 & 0 \checkmark \\
B & n=8 指纹 = 位反序 & $\le$ 0 & 0 \checkmark \\
B & n=16 指纹是置换 & $\le$ 0 & 0 \checkmark \\
B & n=16 指纹 = 位反序 & $\le$ 0 & 0 \checkmark \\
B & n=32 指纹是置换 & $\le$ 0 & 0 \checkmark \\
B & n=32 指纹 = 位反序 & $\le$ 0 & 0 \checkmark \\
A & n=8 的指纹 & [0 4 2 6 1 5 3 7] & $7.352\times 10^{7}$ \checkmark \\
B & n=4 级数 = log2(n) & $\le$ 2.5 & 2 \checkmark \\
B & n=4 逐级残差 & $\le$ 1.42109e-14 & 0 \checkmark \\
B & n=8 级数 = log2(n) & $\le$ 3.5 & 3 \checkmark \\
B & n=8 逐级残差 & $\le$ 2.13163e-14 & 0 \checkmark \\
B & n=16 级数 = log2(n) & $\le$ 4.5 & 4 \checkmark \\
B & n=16 逐级残差 & $\le$ 2.84217e-14 & 0 \checkmark \\
B & n=32 级数 = log2(n) & $\le$ 5.5 & 5 \checkmark \\
B & n=32 逐级残差 & $\le$ 3.55271e-14 & $3.416\times 10^{-16}$ \checkmark \\
B & n=64 级数 = log2(n) & $\le$ 6.5 & 6 \checkmark \\
B & n=64 逐级残差 & $\le$ 4.26326e-14 & $1.364\times 10^{-15}$ \checkmark \\
B & n=128 级数 = log2(n) & $\le$ 7.5 & 7 \checkmark \\
B & n=128 逐级残差 & $\le$ 4.97380e-14 & $5.150\times 10^{-15}$ \checkmark \\
B & n=256 级数 = log2(n) & $\le$ 8.5 & 8 \checkmark \\
B & n=256 逐级残差 & $\le$ 5.68434e-14 & $5.403\times 10^{-15}$ \checkmark \\
B & n=512 级数 = log2(n) & $\le$ 9.5 & 9 \checkmark \\
B & n=512 逐级残差 & $\le$ 6.39488e-14 & $9.684\times 10^{-15}$ \checkmark \\
B & n=1024 级数 = log2(n) & $\le$ 10.5 & 10 \checkmark \\
B & n=1024 逐级残差 & $\le$ 7.10543e-14 & $2.846\times 10^{-14}$ \checkmark \\
B & 逐级残差最坏值 & $\le$ 64$\cdot$log2(\allowbreak{}n)$\cdot$u$\cdot$$\|$x$\|$$\infty$ & $2.846\times 10^{-14}$ \checkmark \\
B & n=4 单位模 & $\le$ 3.55271e-15 & 0 \checkmark \\
B & n=4 半周期反号 & $\le$ 3.55271e-15 & $1.225\times 10^{-16}$ \checkmark \\
B & n=4 共轭对称 & $\le$ 3.55271e-15 & $2.449\times 10^{-16}$ \checkmark \\
B & n=4 逆 = 共轭 & $\le$ 3.55271e-15 & 0 \checkmark \\
B & n=8 单位模 & $\le$ 7.10543e-15 & $2.220\times 10^{-16}$ \checkmark \\
B & n=8 半周期反号 & $\le$ 7.10543e-15 & $1.570\times 10^{-16}$ \checkmark \\
B & n=8 共轭对称 & $\le$ 7.10543e-15 & $3.140\times 10^{-16}$ \checkmark \\
B & n=8 逆 = 共轭 & $\le$ 7.10543e-15 & 0 \checkmark \\
B & n=16 单位模 & $\le$ 1.42109e-14 & $2.220\times 10^{-16}$ \checkmark \\
B & n=16 半周期反号 & $\le$ 1.42109e-14 & $5.467\times 10^{-16}$ \checkmark \\
B & n=16 共轭对称 & $\le$ 1.42109e-14 & $6.497\times 10^{-16}$ \checkmark \\
B & n=16 逆 = 共轭 & $\le$ 1.42109e-14 & 0 \checkmark \\
B & n=32 单位模 & $\le$ 2.84217e-14 & $2.220\times 10^{-16}$ \checkmark \\
B & n=32 半周期反号 & $\le$ 2.84217e-14 & $5.467\times 10^{-16}$ \checkmark \\
B & n=32 共轭对称 & $\le$ 2.84217e-14 & $6.497\times 10^{-16}$ \checkmark \\
B & n=32 逆 = 共轭 & $\le$ 2.84217e-14 & 0 \checkmark \\
B & n=64 单位模 & $\le$ 5.68434e-14 & $2.220\times 10^{-16}$ \checkmark \\
B & n=64 半周期反号 & $\le$ 5.68434e-14 & $5.551\times 10^{-16}$ \checkmark \\
B & n=64 共轭对称 & $\le$ 5.68434e-14 & $7.109\times 10^{-16}$ \checkmark \\
B & n=64 逆 = 共轭 & $\le$ 5.68434e-14 & 0 \checkmark \\
B & n=256 单位模 & $\le$ 2.27374e-13 & $2.220\times 10^{-16}$ \checkmark \\
B & n=256 半周期反号 & $\le$ 2.27374e-13 & $5.979\times 10^{-16}$ \checkmark \\
B & n=256 共轭对称 & $\le$ 2.27374e-13 & $7.109\times 10^{-16}$ \checkmark \\
B & n=256 逆 = 共轭 & $\le$ 2.27374e-13 & 0 \checkmark \\
B & twiddle 恒等式最坏偏差 & $\le$ 8nu & $7.109\times 10^{-16}$ \checkmark \\
B & n=8 rfft 长度 & $\le$ 5 & 5 \checkmark \\
B & n=8 rfft vs fft 前缀 & $\le$ 5.68434e-14 & 0 \checkmark \\
B & n=8 irfft 往返 & $\le$ 2.27374e-13 & $1.112\times 10^{-16}$ \checkmark \\
B & n=9 rfft 长度 & $\le$ 5 & 5 \checkmark \\
B & n=9 rfft vs fft 前缀 & $\le$ 6.39488e-14 & 0 \checkmark \\
B & n=9 irfft 往返 & $\le$ 2.55795e-13 & $2.446\times 10^{-15}$ \checkmark \\
B & n=16 rfft 长度 & $\le$ 9 & 9 \checkmark \\
B & n=16 rfft vs fft 前缀 & $\le$ 1.13687e-13 & 0 \checkmark \\
B & n=16 irfft 往返 & $\le$ 4.54747e-13 & $8.677\times 10^{-17}$ \checkmark \\
B & n=17 rfft 长度 & $\le$ 9 & 9 \checkmark \\
B & n=17 rfft vs fft 前缀 & $\le$ 1.20792e-13 & 0 \checkmark \\
B & n=17 irfft 往返 & $\le$ 4.83169e-13 & $2.060\times 10^{-15}$ \checkmark \\
B & n=32 rfft 长度 & $\le$ 17 & 17 \checkmark \\
B & n=32 rfft vs fft 前缀 & $\le$ 2.27374e-13 & 0 \checkmark \\
B & n=32 irfft 往返 & $\le$ 9.09495e-13 & $2.406\times 10^{-16}$ \checkmark \\
B & n=100 rfft 长度 & $\le$ 51 & 51 \checkmark \\
B & n=100 rfft vs fft 前缀 & $\le$ 7.10543e-13 & 0 \checkmark \\
B & n=100 irfft 往返 & $\le$ 2.84217e-12 & $2.698\times 10^{-14}$ \checkmark \\
B & n=101 rfft 长度 & $\le$ 51 & 51 \checkmark \\
B & n=101 rfft vs fft 前缀 & $\le$ 7.17648e-13 & 0 \checkmark \\
B & n=101 irfft 往返 & $\le$ 2.87059e-12 & $3.259\times 10^{-14}$ \checkmark \\
B & n=128 rfft 长度 & $\le$ 65 & 65 \checkmark \\
B & n=128 rfft vs fft 前缀 & $\le$ 9.09495e-13 & 0 \checkmark \\
B & n=128 irfft 往返 & $\le$ 3.63798e-12 & $4.228\times 10^{-16}$ \checkmark \\
B & rfft/irfft 最坏相对误差 & $\le$ 256nu & $3.259\times 10^{-14}$ \checkmark \\
B & n=16 pad=2 前缀不变 & $\le$ 0 & 0 \checkmark \\
B & n=16 pad=2 尾部全 0 & $\le$ 0 & 0 \checkmark \\
B & n=16 pad=2 原栅格回原值 & $\le$ 9.09495e-13 & 0 \checkmark \\
B & n=16 pad=4 前缀不变 & $\le$ 0 & 0 \checkmark \\
B & n=16 pad=4 尾部全 0 & $\le$ 0 & 0 \checkmark \\
B & n=16 pad=4 原栅格回原值 & $\le$ 1.81899e-12 & 0 \checkmark \\
B & n=16 pad=8 前缀不变 & $\le$ 0 & 0 \checkmark \\
B & n=16 pad=8 尾部全 0 & $\le$ 0 & 0 \checkmark \\
B & n=16 pad=8 原栅格回原值 & $\le$ 3.63798e-12 & 0 \checkmark \\
B & n=24 pad=2 前缀不变 & $\le$ 0 & 0 \checkmark \\
B & n=24 pad=2 尾部全 0 & $\le$ 0 & 0 \checkmark \\
B & n=24 pad=2 原栅格回原值 & $\le$ 1.36424e-12 & $2.237\times 10^{-15}$ \checkmark \\
B & n=24 pad=4 前缀不变 & $\le$ 0 & 0 \checkmark \\
B & n=24 pad=4 尾部全 0 & $\le$ 0 & 0 \checkmark \\
B & n=24 pad=4 原栅格回原值 & $\le$ 2.72848e-12 & $4.378\times 10^{-15}$ \checkmark \\
B & n=24 pad=8 前缀不变 & $\le$ 0 & 0 \checkmark \\
B & n=24 pad=8 尾部全 0 & $\le$ 0 & 0 \checkmark \\
B & n=24 pad=8 原栅格回原值 & $\le$ 5.45697e-12 & $1.862\times 10^{-14}$ \checkmark \\
B & n=32 pad=2 前缀不变 & $\le$ 0 & 0 \checkmark \\
B & n=32 pad=2 尾部全 0 & $\le$ 0 & 0 \checkmark \\
B & n=32 pad=2 原栅格回原值 & $\le$ 1.81899e-12 & 0 \checkmark \\
B & n=32 pad=4 前缀不变 & $\le$ 0 & 0 \checkmark \\
B & n=32 pad=4 尾部全 0 & $\le$ 0 & 0 \checkmark \\
B & n=32 pad=4 原栅格回原值 & $\le$ 3.63798e-12 & 0 \checkmark \\
B & n=32 pad=8 前缀不变 & $\le$ 0 & 0 \checkmark \\
B & n=32 pad=8 尾部全 0 & $\le$ 0 & 0 \checkmark \\
B & n=32 pad=8 原栅格回原值 & $\le$ 7.27596e-12 & 0 \checkmark \\
B & n=100 pad=2 前缀不变 & $\le$ 0 & 0 \checkmark \\
B & n=100 pad=2 尾部全 0 & $\le$ 0 & 0 \checkmark \\
B & n=100 pad=2 原栅格回原值 & $\le$ 5.68434e-12 & $7.244\times 10^{-15}$ \checkmark \\
B & n=100 pad=4 前缀不变 & $\le$ 0 & 0 \checkmark \\
B & n=100 pad=4 尾部全 0 & $\le$ 0 & 0 \checkmark \\
B & n=100 pad=4 原栅格回原值 & $\le$ 1.13687e-11 & $1.571\times 10^{-14}$ \checkmark \\
B & n=100 pad=8 前缀不变 & $\le$ 0 & 0 \checkmark \\
B & n=100 pad=8 尾部全 0 & $\le$ 0 & 0 \checkmark \\
B & n=100 pad=8 原栅格回原值 & $\le$ 2.27374e-11 & $3.315\times 10^{-14}$ \checkmark \\
B & 补零后原栅格回原值的最坏相对差 & X\textasciicircum{}(\allowbreak{}m)\_\{k$\cdot$pad\} =\allowbreak{} X\_k & $3.315\times 10^{-14}$ \checkmark \\
B & 叠加性 & $\le$ 4.54747e-13 & $2.488\times 10^{-16}$ \checkmark \\
B & 齐次性 & $\le$ 4.54747e-13 & $3.926\times 10^{-16}$ \checkmark \\
B & 位移不改幅度谱 & $\le$ 4.54747e-13 & $4.248\times 10^{-16}$ \checkmark \\
B & 位移 = 相位斜坡 & $\le$ 1.81899e-12 & $1.027\times 10^{-14}$ \checkmark \\
B & 位移定理最坏相对误差 & X\_k e\textasciicircum{}\{-2$\pi$ikp/\allowbreak{}N\} & $1.027\times 10^{-14}$ \checkmark \\
B & n=8 栅格长度 & $\le$ 8.5 & 8 \checkmark \\
B & n=8 dt=1 栅格位置 & $\le$ 8.88178e-16 & 0 \checkmark \\
B & n=8 单边谱长度 & $\le$ 5.5 & 5 \checkmark \\
B & n=8 栅格长度 & $\le$ 8.5 & 8 \checkmark \\
B & n=8 dt=0.25 栅格位置 & $\le$ 3.55271e-15 & 0 \checkmark \\
B & n=8 单边谱长度 & $\le$ 5.5 & 5 \checkmark \\
B & n=8 栅格长度 & $\le$ 8.5 & 8 \checkmark \\
B & n=8 dt=0.015625 栅格位置 & $\le$ 5.68434e-14 & 0 \checkmark \\
B & n=8 单边谱长度 & $\le$ 5.5 & 5 \checkmark \\
B & n=9 栅格长度 & $\le$ 9.5 & 9 \checkmark \\
B & n=9 dt=1 栅格位置 & $\le$ 8.88178e-16 & $1.110\times 10^{-16}$ \checkmark \\
B & n=9 单边谱长度 & $\le$ 5.5 & 5 \checkmark \\
B & n=9 栅格长度 & $\le$ 9.5 & 9 \checkmark \\
B & n=9 dt=0.25 栅格位置 & $\le$ 3.55271e-15 & $4.441\times 10^{-16}$ \checkmark \\
B & n=9 单边谱长度 & $\le$ 5.5 & 5 \checkmark \\
B & n=9 栅格长度 & $\le$ 9.5 & 9 \checkmark \\
B & n=9 dt=0.015625 栅格位置 & $\le$ 5.68434e-14 & $7.105\times 10^{-15}$ \checkmark \\
B & n=9 单边谱长度 & $\le$ 5.5 & 5 \checkmark \\
B & n=32 栅格长度 & $\le$ 32.5 & 32 \checkmark \\
B & n=32 dt=1 栅格位置 & $\le$ 8.88178e-16 & 0 \checkmark \\
B & n=32 单边谱长度 & $\le$ 17.5 & 17 \checkmark \\
B & n=32 栅格长度 & $\le$ 32.5 & 32 \checkmark \\
B & n=32 dt=0.25 栅格位置 & $\le$ 3.55271e-15 & 0 \checkmark \\
B & n=32 单边谱长度 & $\le$ 17.5 & 17 \checkmark \\
B & n=32 栅格长度 & $\le$ 32.5 & 32 \checkmark \\
B & n=32 dt=0.015625 栅格位置 & $\le$ 5.68434e-14 & 0 \checkmark \\
B & n=32 单边谱长度 & $\le$ 17.5 & 17 \checkmark \\
B & n=101 栅格长度 & $\le$ 101.5 & 101 \checkmark \\
B & n=101 dt=1 栅格位置 & $\le$ 8.88178e-16 & $1.110\times 10^{-16}$ \checkmark \\
B & n=101 单边谱长度 & $\le$ 51.5 & 51 \checkmark \\
B & n=101 栅格长度 & $\le$ 101.5 & 101 \checkmark \\
B & n=101 dt=0.25 栅格位置 & $\le$ 3.55271e-15 & $4.441\times 10^{-16}$ \checkmark \\
B & n=101 单边谱长度 & $\le$ 51.5 & 51 \checkmark \\
B & n=101 栅格长度 & $\le$ 101.5 & 101 \checkmark \\
B & n=101 dt=0.015625 栅格位置 & $\le$ 5.68434e-14 & $7.105\times 10^{-15}$ \checkmark \\
B & n=101 单边谱长度 & $\le$ 51.5 & 51 \checkmark \\
B & n=256 栅格长度 & $\le$ 256.5 & 256 \checkmark \\
B & n=256 dt=1 栅格位置 & $\le$ 8.88178e-16 & 0 \checkmark \\
B & n=256 单边谱长度 & $\le$ 129.5 & 129 \checkmark \\
B & n=256 栅格长度 & $\le$ 256.5 & 256 \checkmark \\
B & n=256 dt=0.25 栅格位置 & $\le$ 3.55271e-15 & 0 \checkmark \\
B & n=256 单边谱长度 & $\le$ 129.5 & 129 \checkmark \\
B & n=256 栅格长度 & $\le$ 256.5 & 256 \checkmark \\
B & n=256 dt=0.015625 栅格位置 & $\le$ 5.68434e-14 & 0 \checkmark \\
B & n=256 单边谱长度 & $\le$ 129.5 & 129 \checkmark \\
B & 栅格位置的相对偏差 & f\_k =\allowbreak{} k/\allowbreak{}(\allowbreak{}N$\cdot$dt)，相对偏差 $\le$ 64u(\allowbreak{}不是逐位相等) & $1.121\times 10^{-14}$ \checkmark \\
B & n=1 m $\ge$ 2n-1 & $\le$ 1 & 1 \checkmark \\
B & n=1 m 是 2 的幂 & $\le$ 0 & 0 \checkmark \\
B & n=2 m $\ge$ 2n-1 & $\le$ 4 & 3 \checkmark \\
B & n=2 m 是 2 的幂 & $\le$ 0 & 0 \checkmark \\
B & n=3 m $\ge$ 2n-1 & $\le$ 8 & 5 \checkmark \\
B & n=3 m 是 2 的幂 & $\le$ 0 & 0 \checkmark \\
B & n=5 m $\ge$ 2n-1 & $\le$ 16 & 9 \checkmark \\
B & n=5 m 是 2 的幂 & $\le$ 0 & 0 \checkmark \\
B & n=17 m $\ge$ 2n-1 & $\le$ 64 & 33 \checkmark \\
B & n=17 m 是 2 的幂 & $\le$ 0 & 0 \checkmark \\
B & n=100 m $\ge$ 2n-1 & $\le$ 256 & 199 \checkmark \\
B & n=100 m 是 2 的幂 & $\le$ 0 & 0 \checkmark \\
B & n=1000 m $\ge$ 2n-1 & $\le$ 2048 & 1999 \checkmark \\
B & n=1000 m 是 2 的幂 & $\le$ 0 & 0 \checkmark \\
B & n=4096 m $\ge$ 2n-1 & $\le$ 8192 & 8191 \checkmark \\
B & n=4096 m 是 2 的幂 & $\le$ 0 & 0 \checkmark \\
B & n=1 next\_\allowbreak{}power\_\allowbreak{}of\_\allowbreak{}two $\ge$ n & $\le$ 1 & 1 \checkmark \\
B & n=1 next\_\allowbreak{}power\_\allowbreak{}of\_\allowbreak{}two 是 2 的幂 & $\le$ 0 & 0 \checkmark \\
B & n=2 next\_\allowbreak{}power\_\allowbreak{}of\_\allowbreak{}two $\ge$ n & $\le$ 2 & 2 \checkmark \\
B & n=2 next\_\allowbreak{}power\_\allowbreak{}of\_\allowbreak{}two 是 2 的幂 & $\le$ 0 & 0 \checkmark \\
B & n=3 next\_\allowbreak{}power\_\allowbreak{}of\_\allowbreak{}two $\ge$ n & $\le$ 4 & 3 \checkmark \\
B & n=3 next\_\allowbreak{}power\_\allowbreak{}of\_\allowbreak{}two 是 2 的幂 & $\le$ 0 & 0 \checkmark \\
B & n=4 next\_\allowbreak{}power\_\allowbreak{}of\_\allowbreak{}two $\ge$ n & $\le$ 4 & 4 \checkmark \\
B & n=4 next\_\allowbreak{}power\_\allowbreak{}of\_\allowbreak{}two 是 2 的幂 & $\le$ 0 & 0 \checkmark \\
B & n=5 next\_\allowbreak{}power\_\allowbreak{}of\_\allowbreak{}two $\ge$ n & $\le$ 8 & 5 \checkmark \\
B & n=5 next\_\allowbreak{}power\_\allowbreak{}of\_\allowbreak{}two 是 2 的幂 & $\le$ 0 & 0 \checkmark \\
B & n=8 next\_\allowbreak{}power\_\allowbreak{}of\_\allowbreak{}two $\ge$ n & $\le$ 8 & 8 \checkmark \\
B & n=8 next\_\allowbreak{}power\_\allowbreak{}of\_\allowbreak{}two 是 2 的幂 & $\le$ 0 & 0 \checkmark \\
B & n=9 next\_\allowbreak{}power\_\allowbreak{}of\_\allowbreak{}two $\ge$ n & $\le$ 16 & 9 \checkmark \\
B & n=9 next\_\allowbreak{}power\_\allowbreak{}of\_\allowbreak{}two 是 2 的幂 & $\le$ 0 & 0 \checkmark \\
B & n=16 next\_\allowbreak{}power\_\allowbreak{}of\_\allowbreak{}two $\ge$ n & $\le$ 16 & 16 \checkmark \\
B & n=16 next\_\allowbreak{}power\_\allowbreak{}of\_\allowbreak{}two 是 2 的幂 & $\le$ 0 & 0 \checkmark \\
B & n=17 next\_\allowbreak{}power\_\allowbreak{}of\_\allowbreak{}two $\ge$ n & $\le$ 32 & 17 \checkmark \\
B & n=17 next\_\allowbreak{}power\_\allowbreak{}of\_\allowbreak{}two 是 2 的幂 & $\le$ 0 & 0 \checkmark \\
B & n=33 next\_\allowbreak{}power\_\allowbreak{}of\_\allowbreak{}two $\ge$ n & $\le$ 64 & 33 \checkmark \\
B & n=33 next\_\allowbreak{}power\_\allowbreak{}of\_\allowbreak{}two 是 2 的幂 & $\le$ 0 & 0 \checkmark \\
B & n=4096 的 Bluestein/radix-2 相对误差 & $\le$ 64$\cdot$n$\cdot$u & $2.433\times 10^{-12}$ \checkmark \\
B & n=4 plan.n & $\le$ 4.5 & 4 \checkmark \\
B & n=4 有计划名 & $\le$ 0 & 0 \checkmark \\
B & n=8 plan.n & $\le$ 8.5 & 8 \checkmark \\
B & n=8 有计划名 & $\le$ 0 & 0 \checkmark \\
B & n=64 plan.n & $\le$ 64.5 & 64 \checkmark \\
B & n=64 有计划名 & $\le$ 0 & 0 \checkmark \\
B & n=1024 plan.n & $\le$ 1024.5 & 1024 \checkmark \\
B & n=1024 有计划名 & $\le$ 0 & 0 \checkmark \\
B & n=3 有计划名 & $\le$ 0 & 0 \checkmark \\
B & n=5 有计划名 & $\le$ 0 & 0 \checkmark \\
B & n=7 有计划名 & $\le$ 0 & 0 \checkmark \\
B & n=100 有计划名 & $\le$ 0 & 0 \checkmark \\
B & n=127 有计划名 & $\le$ 0 & 0 \checkmark \\
B & 计划不影响结果(seed 0) & $\le$ 1.00000e-10 & $6.405\times 10^{-15}$ \checkmark \\
B & 计划不影响结果(seed 1) & $\le$ 1.00000e-10 & $4.441\times 10^{-15}$ \checkmark \\
B & dft\_\allowbreak{}vs\_\allowbreak{}fft\_\allowbreak{}error(64) & $\le$ 64$\cdot$64$\cdot$u & $6.405\times 10^{-15}$ \checkmark \\
B & n=4 朴素 DIF 末行重排后 = DFT & $\le$ 1.13687e-13 & $4.698\times 10^{-17}$ \checkmark \\
B & n=4 第一级 = 两半的和/差 & $\le$ 3.55271e-15 & 0 \checkmark \\
B & n=8 朴素 DIF 末行重排后 = DFT & $\le$ 2.27374e-13 & $1.045\times 10^{-16}$ \checkmark \\
B & n=8 第一级 = 两半的和/差 & $\le$ 7.10543e-15 & 0 \checkmark \\
B & n=16 朴素 DIF 末行重排后 = DFT & $\le$ 4.54747e-13 & $1.724\times 10^{-16}$ \checkmark \\
B & n=16 第一级 = 两半的和/差 & $\le$ 1.42109e-14 & 0 \checkmark \\
B & n=32 朴素 DIF 末行重排后 = DFT & $\le$ 9.09495e-13 & $1.585\times 10^{-16}$ \checkmark \\
B & n=32 第一级 = 两半的和/差 & $\le$ 2.84217e-14 & 0 \checkmark \\
B & n=64 朴素 DIF 末行重排后 = DFT & $\le$ 1.81899e-12 & $1.941\times 10^{-16}$ \checkmark \\
B & n=64 第一级 = 两半的和/差 & $\le$ 5.68434e-14 & 0 \checkmark \\
B & 朴素 DIF 末行重排后最坏误差 & $\le$ 256nu & $1.941\times 10^{-16}$ \checkmark \\
B & n=4 定义式往返 & $\le$ 2.84217e-14 & $1.349\times 10^{-16}$ \checkmark \\
B & n=7 定义式往返 & $\le$ 4.97380e-14 & $1.930\times 10^{-16}$ \checkmark \\
B & n=16 定义式往返 & $\le$ 1.13687e-13 & $3.829\times 10^{-16}$ \checkmark \\
B & n=33 定义式往返 & $\le$ 2.34479e-13 & $3.395\times 10^{-16}$ \checkmark \\
B & n=64 定义式往返 & $\le$ 4.54747e-13 & $3.791\times 10^{-16}$ \checkmark \\
B & W$\cdot$conj(W)/N = I & $\le$ 1.13687e-13 & $1.193\times 10^{-16}$ \checkmark \\
B & 定义式往返最坏误差 & $\le$ 64nu & $3.829\times 10^{-16}$ \checkmark \\
B & 幅度 0.25 在 k=5 & 0.25 & 0.25 \checkmark \\
B & 峰位 k=5 & 5 & 5 \checkmark \\
B & 幅度 0.25 在 k=17 & 0.25 & 0.25 \checkmark \\
B & 峰位 k=17 & 17 & 17 \checkmark \\
B & 幅度 0.25 在 k=60 & 0.25 & 0.25 \checkmark \\
B & 峰位 k=60 & 60 & 60 \checkmark \\
B & 幅度 0.25 在 k=100 & 0.25 & 0.25 \checkmark \\
B & 峰位 k=100 & 100 & 100 \checkmark \\
B & 幅度 1.0 在 k=5 & 1 & 1 \checkmark \\
B & 峰位 k=5 & 5 & 5 \checkmark \\
B & 幅度 1.0 在 k=17 & 1 & 1 \checkmark \\
B & 峰位 k=17 & 17 & 17 \checkmark \\
B & 幅度 1.0 在 k=60 & 1 & 1 \checkmark \\
B & 峰位 k=60 & 60 & 60 \checkmark \\
B & 幅度 1.0 在 k=100 & 1 & 1 \checkmark \\
B & 峰位 k=100 & 100 & 100 \checkmark \\
B & 幅度 3.5 在 k=5 & 3.5 & 3.5 \checkmark \\
B & 峰位 k=5 & 5 & 5 \checkmark \\
B & 幅度 3.5 在 k=17 & 3.5 & 3.5 \checkmark \\
B & 峰位 k=17 & 17 & 17 \checkmark \\
B & 幅度 3.5 在 k=60 & 3.5 & 3.5 \checkmark \\
B & 峰位 k=60 & 60 & 60 \checkmark \\
B & 幅度 3.5 在 k=100 & 3.5 & 3.5 \checkmark \\
B & 峰位 k=100 & 100 & 100 \checkmark \\
B & 直流读数 = 2.5 & 2.5 & 2.5 \checkmark \\
B & Nyquist 读数 = 1.75 & 1.75 & 1.75 \checkmark \\
B & Nyquist 信号无直流分量 & 0 & 0 \checkmark \\
A & 三种标定场景的最大读数偏差 & 不超过 1e-12(\allowbreak{}绝对) & $1.332\times 10^{-15}$ \checkmark \\
B & n=64 全谱 Parseval & $\le$ 4.54747e-13 & $2.274\times 10^{-16}$ \checkmark \\
B & n=64 单边 Parseval & $\le$ 4.54747e-13 & $2.274\times 10^{-16}$ \checkmark \\
B & n=100 全谱 Parseval & $\le$ 7.10543e-13 & $2.453\times 10^{-16}$ \checkmark \\
B & n=100 单边 Parseval & $\le$ 7.10543e-13 & $2.576\times 10^{-15}$ \checkmark \\
B & n=101 全谱 Parseval & $\le$ 7.17648e-13 & $1.204\times 10^{-16}$ \checkmark \\
B & n=101 单边 Parseval & $\le$ 7.17648e-13 & $1.806\times 10^{-15}$ \checkmark \\
B & n=127 全谱 Parseval & $\le$ 9.02389e-13 & $7.085\times 10^{-16}$ \checkmark \\
B & n=127 单边 Parseval & $\le$ 9.02389e-13 & $1.948\times 10^{-15}$ \checkmark \\
B & n=128 全谱 Parseval & $\le$ 9.09495e-13 & 0 \checkmark \\
B & n=128 单边 Parseval & $\le$ 9.09495e-13 & $1.766\times 10^{-16}$ \checkmark \\
B & n=255 全谱 Parseval & $\le$ 1.81188e-12 & $5.626\times 10^{-16}$ \checkmark \\
B & n=255 单边 Parseval & $\le$ 1.81188e-12 & $3.151\times 10^{-15}$ \checkmark \\
B & n=256 全谱 Parseval & $\le$ 1.81899e-12 & $1.215\times 10^{-16}$ \checkmark \\
B & n=256 单边 Parseval & $\le$ 1.81899e-12 & $2.430\times 10^{-16}$ \checkmark \\
B & n=1000 全谱 Parseval & $\le$ 7.10543e-12 & $2.080\times 10^{-16}$ \checkmark \\
B & n=1000 单边 Parseval & $\le$ 7.10543e-12 & $1.456\times 10^{-15}$ \checkmark \\
A & 单边 df 用错长度时的能量比 & df 用单边谱长度：比值 =\allowbreak{} N/\allowbreak{}(\allowbreak{}N/\allowbreak{}2+1) =\allowbreak{} 1.9804(\allowbreak{}能量被放大近 2 倍) & 1.98039 \checkmark \\
B & N=8,M=5 定理残差 & $\le$ 1.13687e-13 & $2.037\times 10^{-16}$ \checkmark \\
B & N=8,M=5 补零后 = 线性卷积 & $\le$ 9.09495e-13 & $1.538\times 10^{-16}$ \checkmark \\
B & N=16,M=8 定理残差 & $\le$ 2.27374e-13 & $1.898\times 10^{-16}$ \checkmark \\
B & N=16,M=8 补零后 = 线性卷积 & $\le$ 1.81899e-12 & $1.940\times 10^{-16}$ \checkmark \\
B & N=32,M=7 定理残差 & $\le$ 4.54747e-13 & $4.196\times 10^{-16}$ \checkmark \\
B & N=32,M=7 补零后 = 线性卷积 & $\le$ 3.63798e-12 & $3.108\times 10^{-16}$ \checkmark \\
B & N=64,M=13 定理残差 & $\le$ 9.09495e-13 & $2.701\times 10^{-16}$ \checkmark \\
B & N=64,M=13 补零后 = 线性卷积 & $\le$ 7.27596e-12 & $3.346\times 10^{-16}$ \checkmark \\
B & N=12,M=12 定理残差 & $\le$ 1.70530e-13 & $3.965\times 10^{-15}$ \checkmark \\
B & N=12,M=12 补零后 = 线性卷积 & $\le$ 1.81899e-12 & $3.077\times 10^{-16}$ \checkmark \\
A & 循环卷积定理残差 & 不超过 128*N*u & $3.965\times 10^{-15}$ \checkmark \\
B & N=16 差异只出现在 M-1 个位置 & $\le$ 7.5 & 7 \checkmark \\
B & N=16 相对卷绕不超过 1 & $\le$ 1 & 0.421242 \checkmark \\
B & N=24 差异只出现在 M-1 个位置 & $\le$ 7.5 & 7 \checkmark \\
B & N=24 相对卷绕不超过 1 & $\le$ 1 & 0.0809407 \checkmark \\
B & N=32 差异只出现在 M-1 个位置 & $\le$ 7.5 & 7 \checkmark \\
B & N=32 相对卷绕不超过 1 & $\le$ 1 & 0.0242764 \checkmark \\
B & N=64 差异只出现在 M-1 个位置 & $\le$ 7.5 & 7 \checkmark \\
B & N=64 相对卷绕不超过 1 & $\le$ 1 & 0.00274393 \checkmark \\
B & N=128 差异只出现在 M-1 个位置 & $\le$ 7.5 & 7 \checkmark \\
B & N=128 相对卷绕不超过 1 & $\le$ 1 & 0.000863056 \checkmark \\
B & 卷绕误差随 N 单调不增 & $\le$ 0.250781 & 0.000853877 \checkmark \\
B & N=128 时卷绕已经很小但仍非零 & $\le$ 0.01 & 0.000853877 \checkmark \\
A & 补零到 N+M-1 后的最大偏差 & $\le$ 1e-15(\allowbreak{}舍入地板；卷绕的 O(\allowbreak{}1) 项消失) & $1.222\times 10^{-16}$ \checkmark \\
B & rect 峰值旁瓣 vs Harris & $\le$ 0.6 & 0.00145913 \checkmark \\
B & rect -3 dB 主瓣宽 vs Harris & $\le$ 0.05 & 0.00551357 \checkmark \\
B & hann 峰值旁瓣 vs Harris & $\le$ 0.6 & 0.00269145 \checkmark \\
B & hann -3 dB 主瓣宽 vs Harris & $\le$ 0.05 & 0.0017958 \checkmark \\
B & hamming 峰值旁瓣 vs Harris & $\le$ 0.6 & 0.00496879 \checkmark \\
B & hamming -3 dB 主瓣宽 vs Harris & $\le$ 0.05 & 0.000816853 \checkmark \\
B & blackman 峰值旁瓣 vs Harris & $\le$ 0.6 & 0.00115825 \checkmark \\
B & blackman -3 dB 主瓣宽 vs Harris & $\le$ 0.05 & 0.0390711 \checkmark \\
A & 旁瓣次序的实测增量 & 沿 rect < hann < hamming < blackman 的 |旁瓣| 单调增，增量全为正 & 11.2077 \checkmark \\
B & 主瓣宽度：rect 最窄、blackman 最宽 & $\le$ 1.68 & 0.89 \checkmark \\
A & 与 Harris 表的最大偏差 & 不超过 0.6 dB /\allowbreak{} 0.05 格 & 0.00496879 \checkmark \\
B & 整周期：次峰是机器零 & $\le$ 1.00000e-12 & $1.305\times 10^{-16}$ \checkmark \\
B & 偏 0.25 格：次峰约 0.33 & 0.3265 & 0.326483 \checkmark \\
B & 偏 0.5 格：次峰 0.700 & 0.7001 & 0.702359 \checkmark \\
B & Hann 的远处旁瓣比矩形低两个数量级 & $\le$ 5.64468e-04 & 0.000446 \checkmark \\
A & 远离主瓣的旁瓣电平(10 格以外) & Hann 应低至少 20 dB & -42.0461 \checkmark \\
B & rect CG = mean(w) & 1 & 1 \checkmark \\
B & rect 除 CG 后回到 2.0 & 2 & 2 \checkmark \\
B & hann CG = mean(w) & 0.498047 & 0.498047 \checkmark \\
B & hann 除 CG 后回到 2.0 & 2 & 2 \checkmark \\
B & hamming CG = mean(w) & 0.538203 & 0.538203 \checkmark \\
B & hamming 除 CG 后回到 2.0 & 2 & 2 \checkmark \\
B & blackman CG = mean(w) & 0.418359 & 0.418359 \checkmark \\
B & blackman 除 CG 后回到 2.0 & 2 & 2 \checkmark \\
B & Hann 的 CG 接近 0.5 & 0.5 & 0.498047 \checkmark \\
A & 四种窗的相干增益 & rect 1、\allowbreak{}hann 0.5、\allowbreak{}hamming 0.54、\allowbreak{}blackman 0.42 & 0.418359 \checkmark \\
B & rect |斜率| = p & $\le$ 1 & 1 \checkmark \\
B & rect p = 1 & 1 & 1 \checkmark \\
B & hann |斜率| = p & $\le$ 3 & 3 \checkmark \\
B & hann p = 3 & 3 & 3 \checkmark \\
B & hamming |斜率| = p & $\le$ 1 & 1 \checkmark \\
B & hamming p = 1 & 1 & 1 \checkmark \\
B & blackman |斜率| = p & $\le$ 3 & 3 \checkmark \\
B & blackman p = 3 & 3 & 3 \checkmark \\
A & 10 到 1000 的包络比 & 矩形 0.01(\allowbreak{}p=\allowbreak{}1)、\allowbreak{}Hann 1e-6(\allowbreak{}p=\allowbreak{}3) & 0.01 \checkmark \\
B & rect 整数格阈值 = 1.0/T & 1 & 1 \checkmark \\
B & hann 整数格阈值 = 1.0/T & 1 & 1 \checkmark \\
B & hamming 整数格阈值 = 1.0/T & 1 & 1 \checkmark \\
B & blackman 整数格阈值 = 1.0/T & 1 & 1 \checkmark \\
B & 矩形窗在 1.5/T 上分得开 & $\le$ 1 & 1 \checkmark \\
B & Hann 在 1.5/T 上分不开 & $\le$ 0 & 0 \checkmark \\
A & 同一间隔上矩形 vs Hann 的判定 & 矩形可分辨、\allowbreak{}Hann 不可分辨 & 1 \checkmark \\
B & rect 比值不超过 1.8 & $\le$ 1.8 & 1.1236 \checkmark \\
B & rect 比值不小于 0.5 & $\le$ 1.1236 & 0.5 \checkmark \\
B & hann 比值不超过 1.8 & $\le$ 1.8 & 0.694444 \checkmark \\
B & hann 比值不小于 0.5 & $\le$ 0.694444 & 0.5 \checkmark \\
B & hamming 比值不超过 1.8 & $\le$ 1.8 & 0.769231 \checkmark \\
B & hamming 比值不小于 0.5 & $\le$ 0.769231 & 0.5 \checkmark \\
B & blackman 比值不超过 1.8 & $\le$ 1.8 & 0.595238 \checkmark \\
B & blackman 比值不小于 0.5 & $\le$ 0.595238 & 0.5 \checkmark \\
A & 阈值/主瓣宽度比值 & 阈值都贴 1.0/\allowbreak{}T，于是比值恰是 1/\allowbreak{}主瓣宽度的倒数关系 & 0.595238 \checkmark \\
B & n=8,k=0 Goertzel vs fft & $\le$ 3.63798e-12 & $3.155\times 10^{-16}$ \checkmark \\
B & n=8,k=1 Goertzel vs fft & $\le$ 3.63798e-12 & $2.103\times 10^{-16}$ \checkmark \\
B & n=8,k=4 Goertzel vs fft & $\le$ 3.63798e-12 & $6.117\times 10^{-16}$ \checkmark \\
B & n=8,k=5 Goertzel vs fft & $\le$ 3.63798e-12 & $8.479\times 10^{-16}$ \checkmark \\
B & n=8,k=7 Goertzel vs fft & $\le$ 3.63798e-12 & $3.681\times 10^{-16}$ \checkmark \\
B & n=16,k=0 Goertzel vs fft & $\le$ 7.27596e-12 & 0 \checkmark \\
B & n=16,k=1 Goertzel vs fft & $\le$ 7.27596e-12 & $5.903\times 10^{-17}$ \checkmark \\
B & n=16,k=8 Goertzel vs fft & $\le$ 7.27596e-12 & $1.041\times 10^{-15}$ \checkmark \\
B & n=16,k=9 Goertzel vs fft & $\le$ 7.27596e-12 & $4.723\times 10^{-16}$ \checkmark \\
B & n=16,k=15 Goertzel vs fft & $\le$ 7.27596e-12 & $2.094\times 10^{-15}$ \checkmark \\
B & n=64,k=0 Goertzel vs fft & $\le$ 2.91038e-11 & $3.344\times 10^{-15}$ \checkmark \\
B & n=64,k=1 Goertzel vs fft & $\le$ 2.91038e-11 & $6.765\times 10^{-15}$ \checkmark \\
B & n=64,k=32 Goertzel vs fft & $\le$ 2.91038e-11 & $1.591\times 10^{-15}$ \checkmark \\
B & n=64,k=33 Goertzel vs fft & $\le$ 2.91038e-11 & $7.514\times 10^{-15}$ \checkmark \\
B & n=64,k=63 Goertzel vs fft & $\le$ 2.91038e-11 & $6.211\times 10^{-15}$ \checkmark \\
B & n=128,k=0 Goertzel vs fft & $\le$ 5.82077e-11 & $8.650\times 10^{-15}$ \checkmark \\
B & n=128,k=1 Goertzel vs fft & $\le$ 5.82077e-11 & $2.798\times 10^{-15}$ \checkmark \\
B & n=128,k=64 Goertzel vs fft & $\le$ 5.82077e-11 & $5.127\times 10^{-14}$ \checkmark \\
B & n=128,k=65 Goertzel vs fft & $\le$ 5.82077e-11 & $5.338\times 10^{-15}$ \checkmark \\
B & n=128,k=127 Goertzel vs fft & $\le$ 5.82077e-11 & $2.492\times 10^{-15}$ \checkmark \\
B & n=17,k=0 Goertzel vs fft & $\le$ 7.73070e-12 & $1.402\times 10^{-16}$ \checkmark \\
B & n=17,k=1 Goertzel vs fft & $\le$ 7.73070e-12 & $1.147\times 10^{-15}$ \checkmark \\
B & n=17,k=8 Goertzel vs fft & $\le$ 7.73070e-12 & $8.321\times 10^{-16}$ \checkmark \\
B & n=17,k=9 Goertzel vs fft & $\le$ 7.73070e-12 & $1.478\times 10^{-15}$ \checkmark \\
B & n=17,k=16 Goertzel vs fft & $\le$ 7.73070e-12 & $2.342\times 10^{-15}$ \checkmark \\
A & 收尾式写反时的结果 & 正确值 4(\allowbreak{}等于 fft 的第 2 项) & 4 \checkmark \\
B & n=32,dt=1 时域能量 & 28.0359 & 28.0359 \checkmark \\
B & n=32,dt=0.01 时域能量 & 0.280359 & 0.280359 \checkmark \\
B & n=64,dt=1 时域能量 & 83.2563 & 83.2563 \checkmark \\
B & n=64,dt=0.01 时域能量 & 0.832563 & 0.832563 \checkmark \\
B & n=101,dt=1 时域能量 & 76.0184 & 76.0184 \checkmark \\
B & n=101,dt=0.01 时域能量 & 0.760184 & 0.760184 \checkmark \\
B & n=256,dt=1 时域能量 & 272.541 & 272.541 \checkmark \\
B & n=256,dt=0.01 时域能量 & 2.72541 & 2.72541 \checkmark \\
B & 时域能量定义式的最大偏差 & 逐位相等 & 0 \checkmark \\
B & 直接读格：误差约半格 & $\le$ 0.6 & 0.5 \checkmark \\
B & 抛物线插值：误差不超过 5\% 栅格 & $\le$ 0.05 & 0.0422127 \checkmark \\
B & 插值定位与栅格间距之比 & 远小于 1 & 0.0422127 \checkmark \\
B & 3 Hz 与 17 Hz 在 fs=20 下的样本逐位相同 & $\le$ 1.00000e-12 & $4.763\times 10^{-14}$ \checkmark \\
B & 3.0 Hz 的谱峰在 3 Hz & 3 & 3 \checkmark \\
B & 3.0 Hz 的幅度仍是 1.0 & 1 & 1 \checkmark \\
B & 17.0 Hz 的谱峰在 3 Hz & 3 & 3 \checkmark \\
B & 17.0 Hz 的幅度仍是 1.0 & 1 & 1 \checkmark \\
B & alias\_\allowbreak{}frequency(17,20)=3 & 3 & 3 \checkmark \\
B & alias\_\allowbreak{}frequency(12,20)=8(12 折到 8，不是 3) & 8 & 8 \checkmark \\
B & alias\_\allowbreak{}frequency(7,10)=3 & 3 & 3 \checkmark \\
B & alias\_\allowbreak{}index(7,10)=3 & $\le$ 3 & 3 \checkmark \\
B & 3 Hz 在 20 Hz 以内的候选真频率个数 & 3、\allowbreak{}17(\allowbreak{}差一个 fs=\allowbreak{}20) & 2 \checkmark \\
B & 20log10 换算 & 0 & $8.672\times 10^{-8}$ \checkmark \\
B & 0 映射到地板 -120 dB & $\le$ 1.00000e-09 & 0 \checkmark \\
B & 地板后无 NaN/Inf & $\le$ 0 & 0 \checkmark \\
B & 全零信号映射到全地板 -60 dB & $\le$ 1.00000e-09 & 0 \checkmark \\
B & 0 的 dB 值(带地板) & -120 dB & -120 \checkmark \\
B & 栅格不随位移改变 & $\le$ 0 & 0 \checkmark \\
B & 位移不改幅度谱 & $\le$ 4.54747e-13 & $2.131\times 10^{-16}$ \checkmark \\
B & 相位斜坡 = -2*pi*k*p/N & $\le$ 1.00000e-09 & $2.266\times 10^{-15}$ \checkmark \\
B & 相位斜坡的最大偏差 & 不超过 1e-9 & $2.266\times 10^{-15}$ \checkmark \\
B & 循环差分 = 定义式(逐位相等) & $\le$ 0 & 0 \checkmark \\
B & 差分的谱乘子 = e\textasciicircum{}\{i w\} - 1 & $\le$ 3.63798e-12 & $1.136\times 10^{-15}$ \checkmark \\
B & 差分 vs 解析导数的波形偏差 = 2sin(w/4)$\cdot$2sin(w/2)/w & 0.146598 & 0.146908 \checkmark \\
B & i*w 对纯正弦给出精确导数 & $\le$ 1.00000e-13 & $8.858\times 10^{-15}$ \checkmark \\
B & 谱微分与一阶差分是两个不同的乘子 & i*w 给出精确导数；e\textasciicircum{}\{iw\}-1 给出精确一阶差分，两者差 2sin(\allowbreak{}w/\allowbreak{}2)/\allowbreak{}w 的幅度与 w/\allowbreak{}2 的相位 & 0.146908 \checkmark \\
B & 不等长：fft\_\allowbreak{}convolve 等于短序列补零后的循环卷积 & $\le$ 2.27374e-13 & $1.934\times 10^{-16}$ \checkmark \\
B & 等长时两条路径一致 & $\le$ 2.27374e-13 & $3.780\times 10^{-16}$ \checkmark \\
B & circular\_\allowbreak{}convolution 拒绝不等长输入 & 抛 ValueError，而不是悄悄补零 & 1 \checkmark \\
B & rect 能量 = N & 1024 & 1024 \checkmark \\
B & hann 能量约为 3N/8 & 384 & 383.625 \checkmark \\
B & blackman 能量约为 0.3045N & 0.3043 & 0.304303 \checkmark \\
B & Hann/rect 能量比 & 约 3/\allowbreak{}8 =\allowbreak{} 0.375 & 0.374634 \checkmark \\
B & Rayleigh = 1/T & 1 & 1 \checkmark \\
B & 栅格 = 1/(N*dt) & 1.5625 & 1.5625 \checkmark \\
B & pad=1 原栅格逐位回到原值 & $\le$ 0 & 0 \checkmark \\
B & pad=1 栅格 = 1/(N*pad*dt) & 1.5625 & 1.5625 \checkmark \\
B & pad=2 原栅格逐位回到原值 & $\le$ 0 & 0 \checkmark \\
B & pad=2 栅格 = 1/(N*pad*dt) & 0.78125 & 0.78125 \checkmark \\
B & pad=4 原栅格逐位回到原值 & $\le$ 0 & 0 \checkmark \\
B & pad=4 栅格 = 1/(N*pad*dt) & 0.390625 & 0.390625 \checkmark \\
B & pad=8 原栅格逐位回到原值 & $\le$ 0 & 0 \checkmark \\
B & pad=8 栅格 = 1/(N*pad*dt) & 0.195312 & 0.195312 \checkmark \\
B & pad=16 原栅格逐位回到原值 & $\le$ 0 & 0 \checkmark \\
B & pad=16 栅格 = 1/(N*pad*dt) & 0.0976562 & 0.0976562 \checkmark \\
B & 峰值不随 pad 改变 & 1 & 1 \checkmark \\
B & 主瓣半高全宽随 pad 的增量 & 读数单调下降(\allowbreak{}量化误差在消失)，但不收敛到 0 & -0.390625 \checkmark \\
B & 收敛值不是 0：主瓣宽度有正的下限 & $\le$ 1.5 & 0.976562 \checkmark \\
B & T=0.5 Rayleigh = 1/T & 2 & 2 \checkmark \\
B & T=2.0 Rayleigh = 1/T & 0.5 & 0.5 \checkmark \\
B & T=4.0 Rayleigh = 1/T & 0.25 & 0.25 \checkmark \\
\end{longtable}


\section{它会怎样骗你（本章版）}

\begin{enumerate}
  \item \textbf{把补零当成提高分辨率}：补零后谱线从 $33$ 个点变成 $257$ 个点，
        但原栅格上的读数逐位不变（$3.3\times10^{-14}$ 即机器零），
        主瓣半高全宽读数只从 $1.5625$ Hz 收敛到 $0.9766$ Hz。
        要分辨两个音必须加长 $T$，不是加长 $N_{\text{pad}}$。
  \item \textbf{只用 2 的幂做对拍}：长度取 $256$ 时会漏掉整条 Bluestein 分支。
        把长度扩到 $\{1,2,3,4,5,7,8,11,13,16,17,31,32,96,127,128,200,1000\}$ 后
        最坏相对误差 $1.46\times10^{-13}$（$n=1000$），素数长度 $127$ 与 $3$ 都覆盖到了。
  \item \textbf{忘记 $1/N$ 却"看着对"}：形状完全正确，只是幅度大了 $N$ 倍，
        取相对峰值 dB 时这个错误会被彻底吃掉。实测 $n=1000$ 时放大 $1000.0$ 倍，
        而 $n=1$ 时放大 $1$ 倍——所以往返测试的长度集合必须包含 $n=1$，
        把这个"测不出来"的长度本身记录在案。
  \item \textbf{用冲激去测位反序}：冲激的谱是常数，对任何置换都不敏感。
        改用复指数后 $n=8$ 的指纹逐位等于 $[0,4,2,6,1,5,3,7]$；
        不重排时 DIF 末行与 DFT 的相对误差是 $1.071$——这个 $O(1)$ 是下标约定，
        不是精度。
  \item \textbf{把"两份实现一致"当成"算法正确"}：逐级残差 $0$（$n=4,8$）到
        $2.85\times10^{-14}$（$n=1024$）量的是向量化与朴素实现是否一致；
        写错两遍也是 $0$。数学正确性要靠与 $O(n^2)$ 定义式对拍（$1.46\times10^{-13}$）。
  \item \textbf{Goertzel 的共轭方向猜错}：收尾式写反（$e^{i\omega}\to e^{-i\omega}$）时，
        冲激输入、$k=0$、$k=N/2$ 全部通过，只有随机信号 $+$ 一般 $k$ 才暴露（相对差 $1.29$）。
        本章的测试同时含随机信号与 $k=1$、$k=N/2\pm1$ 这些"一般"位置。
  \item \textbf{用单边谱的长度算 $\Delta f$}：$n=101$ 时能量被放大 $1.9804$ 倍，
        Parseval 相对误差 $0.980$。这个数"不像粗心"（不是 $1$、不是 $10$），
        最容易被归类为"数值误差"。
  \item \textbf{只看旁瓣、不看代价}：Blackman 的 $-58.11$ dB 比矩形好 $45$ dB，
        但主瓣宽 $0.884\to1.641$ 格、能量 $N\to0.3043N$、相干增益 $1.0000\to0.4184$。
        少报任何一个，选窗的结论都会反过来。
  \item \textbf{把汉明与 Hann 互换}：汉明第一旁瓣 $-42.7$ dB 比 Hann 低 $11$ dB，
        但它是 $p=1$（滚降 $-6$ dB/倍频程），Hann 是 $p=3$（$-18$ dB/倍频程）。
        实测 $10\to1000$ 的包络比：矩形 $1.00\times10^{-2}$、Hann $1.00\times10^{-6}$。
  \item \textbf{把泄漏怪到矩形窗头上}：整周期采样时次峰是机器零（$0.000$），
        偏 $0.25$ 格就到 $0.3265$（$-9.72$ dB）。泄漏的原因是"信号周期与 $T$ 不成整数比"。
  \item \textbf{用二分法找分辨率阈值}：谓词对间隔不单调，半格间隔会把两音并进一个主瓣，
        二分因此收敛到 $1.73/T$ 这种"看着挺像真值"的假穿越。
        本章改用栅格扫描，并明确报出"阈值在整数格上量得"。
  \item \textbf{用 $|$FFT$|$ 去量窗的 $-3$ dB 主瓣宽度}：窗长小时栅格太粗，
        宽度会被量化到整格（rect 会报 $0.44$ 格而不是 $0.884$ 格）。
        本章改用解析窗谱（与 Harris 表最大差 $0.0050$ dB）。
  \item \textbf{把卷绕误差当成"$N$ 不够大"的小问题}：$(N,M)=(16,8)$ 时它 $=0.251$
        （相对 $42.1\%$），且只出现在前 $7=M-1$ 项上（正好在头部）；
        $N$ 翻到 $64$ 仍有 $2.63\times10^{-3}$。它不会自己变成 $0$，
        必须补零到 $N_a+N_b-1$（之后是 $1.22\times10^{-16}$）。
  \item \textbf{拿两类不同的曲线相减当"误差"}：泄漏页早先拿"加窗信号的谱"减"窗的谱"，
        两条曲线不同类，$rel\_error$ 在 $offset=0$ 时恒为 $1.0$、$offset=0.25$ 时是 $0.784$，
        与实现精度无关。改成同类对同类（$|$FFT$|$ vs 解析窗谱）后是 $0.0055\sim0.0132$，
        且与 $offset$ 无关。
  \item \textbf{把理论的旁瓣线画在"形状对、位置错"的地方}：包络的常数项决定它画在哪。
        早先把自变量写成 $f\cdot n$（$n$ 是点数不是时长），矩形窗的理论线被压低
        $64$ 倍 $=36$ dB，而且它"总在实测曲线下方"，看起来还挺保守。
        正确自变量是 $f\cdot T$，锚点取窗自己的最高旁瓣。
  \item \textbf{把"测到 $3$ Hz"当成"信号是 $3$ Hz"}：$f_s=20$ 时 $3$ Hz 与 $17$ Hz
        的样本逐位相同，幅度都读 $1.0$；$20$ Hz 以内 $3$ Hz 的候选真频率至少 $2$ 个。
  \item \textbf{补零之后振幅整整齐齐掉一半}：$2/N$ 里的 $N$ 必须是原始长度。
        用补零后的长度归一化时每个峰掉 $p$ 倍（$p=2$ 时 $0.5$），但谱形完全正常。
        判据是"补零倍数改变时峰值不动"。
  \item \textbf{把断言里的"理论值"当成抄来的常数}：本章有四处写错过，全部留在
        \texttt{pitfalls.md} 里。例如单边谱长度算 $\Delta f$ 我曾写成"能量只剩 $96.5\%$"，
        实测是放大 $1.9804$ 倍（方向都反了）；"波形偏差 $=2\sin(\omega/2)$"实测是
        $0.1469$ 而该式给 $0.2935$（差一倍，正确式是 $2\sin(\omega/4)\cdot2\sin(\omega/2)/\omega$）；
        "补零后误差逐位为 $0$"实测是 $1.11\times10^{-16}$。理论值必须有来源
        （推导/独立实现/文献表），判据要写成带容差的等式，少用"够大就过"的单向判据。
  \item \textbf{把 $i\omega$ 与 $e^{i\omega}-1$ 混用}：在 $\omega=2\pi\cdot3/64$ 处两者模只差
        $0.36\%$，相位差 $\omega/2$，作用在纯正弦上是 $14.69\%$ 的波形偏差。
        $i\omega$ 是精确的谱微分，$e^{i\omega}-1$ 是精确的一阶差分；两个都对，
        但代表不同算子。
\end{enumerate}
